\documentclass[a4paper,11pt]{memoir}

\usepackage[utf8]{inputenc}
\usepackage[T1]{fontenc}
\usepackage[svgnames]{xcolor}
\usepackage{fouriernc}
\usepackage[frenchb]{babel}
% TODO \usepackage[french]{babelbib}
\usepackage{mathpartir}
\usepackage{syntax}
\usepackage{amssymb}
\usepackage{tikz}
\usepackage{amsmath}
\usepackage{float}
\usepackage{url}
\usepackage{xspace}
\usepackage{graphicx}
\usepackage[shortlabels]{enumitem}
\usepackage[backref=page]{hyperref}
\usepackage{links}
\usepackage[a4paper]{geometry}
\usepackage{amsthm}
\usepackage{stmaryrd}
\usepackage[normalem]{ulem}
\usepackage{wasysym}
\usepackage{multirow}
\usepackage{thmtools}
\usepackage{gramster}
\usepackage{switch}
\usepackage{pelagiccod}

\newcommand{\irule}[3]{ \inferrule*[right=(#1)]{#2}{#3} }

% Fix '::' -> ': :' when using babel
\newenvironment{grammarf}{\shorthandoff{:}\begin{grammar}}{\end{grammar}\shorthandon{:}}

\newcommand{\gramisa}[0]{::=~}
\newcommand{\gramor}[0]{|\hspace{3mm}}

\newcommand{\gramlr}[2]{%
  \begin{minipage}{0.2\textwidth}\raggedleft #1\end{minipage}
  \begin{minipage}{0.8\textwidth}#2\end{minipage}
}

% Types
\newcommand{\tInt}{\textsc{Int}\xspace}
\newcommand{\tFloat}{\textsc{Float}\xspace}
\newcommand{\tVoid}{\textsc{Void}\xspace}

\newcommand{\tFV}[1]{\textrm{FreeVars}(#1)}

\def\opxor{\mathbin{\char`\^}}

% Constant
\newcommand{\cNil}{\textsc{Nil}}

% Program & language names
\newcommand{\ctonewspeak}{\textsc{c2newspeak}}
\newcommand{\cpp}{\textsc{cpp}}
\newcommand{\newspeak}{\textsc{Newspeak}}

% Newspeak
\newcommand{\npkDecl}[2]{\textsc{Decl}~(#1) \{ #2 \}}
\newcommand{\npkIf}[3]{\textsc{If}(#1) \{ #2 \} \textsc{Else} \{ #3 \}}
\newcommand{\npkDoWith}[2]{\textsc{Do}\{ #1 \}\textsc{With}~#2:}
\newcommand{\npkWhile}[1]{\textsc{While}(1)\{#1\}}
\newcommand{\npkGoto}[1]{\textsc{Goto}~#1}
\newcommand{\npkReturn}[1]{\textsc{Return}~#1}

% Continuations
\def\kPass{\textsc{Kpass}}
\def\kSeq#1#2{\textsc{Kseq}(#1, #2)}
\def\kBlock#1#2{\textsc{Kblock}(#1, #2)}
\def\kGoto#1{\textsc{Kgoto}(#1)}
\def\kCall#1#2#3#4{\textsc{Kcall}(#1, #2, #3, #4)}
\def\kRet#1{\textsc{Kret}(#1)}

\newcommand{\insertcode}[1]{\input{gen/#1.pyg.tex}}

% \chapter* + add to toc
%\newcommand{\chapterstar}[1]{\chapter*{#1} \addcontentsline{toc}{chapter}{#1}}

% Register names
\def\defreg#1{\expandafter\def\csname#1\endcsname{\textsc{#1}\xspace}}

\defreg{eax}
\defreg{ebx}
\defreg{ecx}
\defreg{edx}
\defreg{esi}
\defreg{edi}
\defreg{ebp}
\defreg{esp}
\defreg{eip}
\defreg{eflags}

\def\crtrois{\textsc{CR3}}

\newcommand{\ring}[1]{\emph{ring}~#1}

\newcommand{\asminstr}[1]{\texttt{#1}}

\newcommand{\wip}[0]{%
\begin{center}
    \includegraphics[width=5cm]{gen/travaux.pdf}
\end{center}
}

\newcommand{\listcons}[0]{%
\texttt{(::)}%
}

\newcommand{\hexa}[1]{%
\texttt{0x#1}%
}

% Debug tikz bounding box
\def\drawbb{\draw[thick] (current bounding box.south west) rectangle (current bounding box.north east);}

% comes from the memoir manual {{{

\newcommand*{\addtotextheight}[1]{%
\twocolumn
\addtolength{\textheight}{#1}%
\onecolumn}

%% for main body, bottom of text at 1in, footer below
%% top of header at 1in, first text line double spaced
%% below base of header
\newlength{\linespace}
\setlength{\linespace}{\baselineskip} % current equivalent of \onelineskip
\setlength{\headheight}{\onelineskip}
\setlength{\headsep}{\linespace}
\addtolength{\headsep}{-\topskip}
\setlength{\uppermargin}{1in}
\addtolength{\uppermargin}{\headheight}
\addtolength{\uppermargin}{\headsep}

%% and for the bottom
\setlength{\lowermargin}{1in}
\setlength{\textheight}{\paperheight}

\addtolength{\textheight}{-\uppermargin}
\addtolength{\textheight}{-\lowermargin}
%% footnote settings
\setlength{\footskip}{\onelineskip}
\setlength{\footnotesep}{\onelineskip}


%% the fiddle lengths (..ta.. for title/approval page, others for prelims)
\newlength{\toptafiddle} \setlength{\toptafiddle}{2\linespace}
\newlength{\bottafiddle} \setlength{\bottafiddle}{0pt}
\newlength{\topfiddle}
\setlength{\topfiddle}{\toptafiddle}
\newlength{\botfiddle}
\setlength{\botfiddle}{\onelineskip}

%%% general macro for Abstract, etc., headings
\newcommand*{\pretoctitle}[1]{{\clearpage\centering \vspace*{-\topfiddle}#1\par}}

%%% Start the ACKNOWLEDGEMENTS
\newcommand{\asuacknowledgements}{\pretoctitle{ACKNOWLEDGEMENTS}}

%%% Start the ABSTRACT
\newcommand{\asuabstract}{ \addtotextheight{-\botfiddle} \pretoctitle{ABSTRACT}}

\addtodef{\mainmatter}{\addtotextheight{\botfiddle}{}}

%% make it easy to center any dedication
\newcommand{\asudedication}[1]{%
{\clearpage\mbox{}\vfill\centering #1 \par\vfill\clearpage}}

% }}} /memoir

\DeclareUnicodeCharacter{0393}{\Gamma}        % Γ
\DeclareUnicodeCharacter{03B5}{\epsilon}      % ε
\DeclareUnicodeCharacter{03C3}{\sigma}        % σ
\DeclareUnicodeCharacter{2190}{\leftarrow}    % ←
\DeclareUnicodeCharacter{2191}{\uparrow}      % ↑
\DeclareUnicodeCharacter{2192}{\rightarrow}   % →
\DeclareUnicodeCharacter{21D2}{\Rightarrow}   % ⇒
\DeclareUnicodeCharacter{2205}{\emptyset}     % ∅
\DeclareUnicodeCharacter{2208}{\in}           % ∈
\DeclareUnicodeCharacter{2260}{\neq}          % ≠
\DeclareUnicodeCharacter{2295}{\oplus}        % ⊕
\DeclareUnicodeCharacter{22A2}{\vdash}        % ⊢


\input{chapstyle.tex}

\usetikzlibrary{calc}
\usetikzlibrary{shapes}
\usetikzlibrary{arrows}
\usetikzlibrary{decorations.pathreplacing}

\tikzstyle{bigbrace} = [decorate,decoration={brace,amplitude=3mm}]


\usepackage{fancyvrb}
\input{gen/pygments-style.tex}

\floatstyle{boxed}
\restylefloat{figure}

\setlist[itemize]{label=$\bullet$}

\isopage[12]
\checkandfixthelayout

\maxtocdepth{subsection}    % What appears in toc
\maxsecnumdepth{subsection} % What get numbered

%\theoremstyle{definition}
\newtheorem{definition}{Définition}

\def\iffull#1{\if#1\relax\else}
\title{ \bfseries\LARGE Analyse statique de logiciel système par typage statique fort\\%
      {\mdseries\Large Application au noyau Linux}}
\author{Étienne Millon}

\begin{document}


\frontmatter

\maketitle

\cleartorecto

%\asuabstract

%Abstract ici.

%\asuacknowledgements

%Remerciements

%\asudedication{Cassedédi}

\tableofcontents

\mainmatter

\chapter{Introduction}
Communication,
audiovisuel,
transports,
médecine :
tout ces domaines se sont transformés dans les dernières décennies,
en particulier grâce à la révolution numérique.
En effet le plus petit appareil électrique contient maintenant des composants
matériels programmables.

% TODO faire une intro "la vie est numérique" avec du Berry @ CdF

En 2013, on pense bien sûr aux téléphones portables dont la fonctionnalité et la
complexité les rapprochent des ordinateurs de bureau. Par exemple, le système
d'exploitation Android de Google est basé sur le noyau Linux, destiné à la base
aux micro-ordinateurs.

Le noyau d'un système d'exploitation est chargé de faire l'intermédiaire entre
le matériel (processeur, mémoire, périphériques, …) et les applications
exécutées sur celui-ci (par exemple un navigateur web, une calculatrice ou un
carnet d'adresses). Il doit aussi garantir la sécurité et l'isolation de
celles-ci.

En tant qu'intermédiaire de confiance, le noyau a un certain nombre de
responsabilités et est le seul à avoir accès à certaines informations sensibles.
Il est capital de s'assurer qu'il est bien le seul à pouvoir y accéder. En
particulier, il faut pouvoir vérifier que les requêtes faites par l'utilisateur
au noyau ne peuvent pas détourner ce dernier et lui faire fuiter des
informations confidentielles.

Le problème est que comme tous les logiciels, les noyaux de système
d'exploitation sont écrits par des humains, qui ne sont pas parfaits. Loin de
là: on estime qu'avant relecture, 1000 lignes de code contiennent entre 5 et 100
erreurs de programmation en moyenne.
% TODO sourcer
Les activités de relecture et de
débogage ont beau prendre la majeure partie du temps de développement, il est
facile de laisser passer des défauts de programmation.

Une technique efficace est de réaliser des tests, c'est-à-dire exécuter le code
du programme sous un environnement contrôlé. On peut alors détecter des
comportements non désirés. Mais même avec des tests les plus exhaustifs il n'est
pas possible de couvrir tous les cas d'utilisation.

Parmi les techniques communes pour éliminer les ``bugs'', on peut ajouter à un
langage de programmation un \emph{système de types statiques}.

Celui-ci permet d'associer à chaque fragment d'un programme, une étiquette
décrivant quel genre de valeur sera produite par son évaluation. Cette analyse,
faite à la compilation, permet de rejeter les programmes incorrects du point de
vue de cette catégorisation.

Le but de cet thèse est de montrer que le typage statique peut être utilisé pour
détecter et interdire certains bugs présents dans les noyaux de systèmes
d'exploitation.

\section{Rôle d'un système d'exploitation}

On décrit ici la fonction générale d'un système d'exploitation. Ce sujet est
détaillé en profondeur, y compris avec des cas d'étude, dans~\cite{tanenbaum}.

Un ordinateur est constitué de nombreux composants matériels : microprocesseur,
mémoire, et divers périphériques. Pourtant, au niveau de l'utilisateur, des
dizaines de logiciels permettent d'effectuer toutes sortes de calculs et de
communications. Le système d'exploitation permet de faire l'interface entre ces
niveaux d'abstraction.

Au cours de l'histoire des systèmes informatiques, la manière de les programmer
a beaucoup évolué. Au départ, les programmeurs avaient accès au matériel dans
son intégralité : toute la mémoire pouvait être accédée, toutes les instructions
pouvaient être utilisées.

Néanmoins c'est un peu restrictif, puisque cela ne permet qu'à une personne
d'interagir avec le système. Dans la seconde moitié des années 60, sont apparus
les premiers systèmes ``à temps partagé'', permettant à plusieurs utilisateurs
de travailler en même temps.

Permettre l'exécution de plusieurs programmes en même temps est une idée
révolutionnaire, mais elle n'est pas sans difficultés techniques : en effet les
ressources de la machine doivent être aussi partagées entre les utilisateurs et
les programmes. Par exemple, plusieurs programmes vont utiliser le processeur
les uns à la suite des autres ; et chaque programme aura à sa disposition une
partie de la mémoire principale, ou du disque dur.

Si deux programmes (ou plus) s'exécutent de manière concurrente sur le même
matériel, il faut s'assurer que l'un ne puisse pas écrire dans la mémoire de
l'autre, ou que les deux n'utilisent pas la carte réseau en même temps. Ce sont
des rôles du système d'exploitation.

Ainsi, au lieu d'accéder directement au matériel via des instructions de bas
niveau, les programmes communiquent avec le noyau, qui centralise donc les
appels au matériel, et abstrait certaines opérations.

Par exemple, comparons ce qui se passe concrètement lors de la copie de données
depuis un cédérom ou une mémorette USB.

\begin{itemize}

  \item Dans le cas du cédérom, il faut interroger le bus SATA, interroger le
    lecteur sur la présence d'un disque dans le lecteur, activer le moteur,
    calculer le numéro de trame des données sur le disque, demander la lecture,
    puis déclencher une
    copie de la mémoire.

  \item Avec une mémorette USB, il faut interroger le bus USB, rechercher le bon
    numéro de périphérique, le bon numéro de canal dans celui-ci, lui appliquer
    une commande de lecture au bon numéro de bloc, puis copier la mémoire.

\end{itemize}

Ces deux opérations, bien qu'elles aient le même but (copier de la mémoire
depuis un périphérique amovible), ne sont pas effectuées de la même manière.
C'est pourquoi le système d'exploitation fournit les notions de fichier,
lecteur, etc : le programmeur n'a plus qu'à utiliser des commandes de haut
niveau (``monter un lecteur'', ``ouvrir un fichier'', ``lire dans un fichier'')
et selon le type de lecteur, le système d'exploitation effectuera les actions
appropriées.

En résumé, un système d'exploitation est l'intermédiaire entre le logiciel et
le matériel, et en particulier assure les rôles suivants :

% TODO À affiner ou supprimer

\paragraph{Gestion des processus}

Un système d'exploitation peut permettre d'exécuter plusieurs programmes à la
fois. Il faut alors orchestrer ces différents processus et les séparer en terme
de temps et de ressources partagées.

\paragraph{Gestion de la mémoire}

Chaque processus, en plus du noyau, doit disposer d'un espace mémoire différent.
C'est-à-dire qu'un processus ne doit pas pouvoir interférer avec un autre.

\paragraph{Gestion des fichiers}

Les processus peuvent accéder à une hiérarchie de fichiers, indépendamment de la
manière d'y accéder.

\paragraph{Gestion des périphériques}

Le noyau étant le seul code ayant des privilèges, c'est lui qui doit communiquer
avec les périphériques matériels.

\paragraph{Abstractions}

Le noyau fournit aux programmes une interface unifiée, permettant de stocker des
informations de la même manière sur un disque dur ou une clef USB (alors que
l'accès se déroulera de manière très différente en pratique).

\section{Noyaux de systèmes d'exploitation}

% TODO [E] ??? privilèges ?

Les programmes qui composent un système d'exploitation sont exécutés avec un
niveau de privilège différent, correspondant à un jeu d'instructions plus ou
moins important. Cette vérification est faite directement pas le processeur,
d'une manière matérielle. En plus de cette distinction, les programmes ont une
vue différente de la mémoire. Si un programme utilisateur avait accès aux
structures de données internes du noyau, il pourrait obtenir des informations
auquel il n'a pas normalement accès, ou obtenir des privilège supplémentaires
(par exemple en obtenant les privilèges de l'administrateur sur la machine).

C'est pourquoi les noyaux stockent leur code et données dans un endroit en
mémoire qui n'est pas accessible aux programmes utilisateur. Sur un système
Intel 32 bits, cela correspond aux adresses supérieures à \texttt{0xc0000000},
ou 3 Gio.

% TODO trad

Because user-mode programs cannot run low-level instructions, they cannot do
anything interesting on their own. They have to use the kernel's facilities to
perform input/output or manipulate the operating system's abstractions (such as
creating a new process). To this end, they can communicate with the kernel
via the narrow interface of \emph{system calls}.

For example, every time a user process calls the \texttt{read()} function, an
architecture-specific mechanism invokes the kernel-mode function
\texttt{sys\_read()}. The user passes a file descriptor, a pointer and an amount
of bytes to fill and the kernel will actually perform input/output and copy the
resulting bytes on the user's buffer.

But if the user supplies a pointer whose value is in the kernel's reserved area,
this zone will be overwritten, potentially leading to a security breach. In a
sense, the kernel has been abused because it accessed the memory with its own
privileges instead of the originator's. This is known as the \emph{confused
deputy problem}\cite{hardy88confused}.

When implementing a system call, it is thus necessary to forbid direct
dereference of a pointer whose value can be controlled by the user. In case of a
system call with a pointer argument, one has to dynamically verify that the
address provided by the user lies within the process's memory space.

% TODO [E] est-ce ici que tu veux le placer ? rien ne t'empêche de le dire
% plusieurs fois

The object of this work is to detect the places where it is necessary to insert
dynamic checks. To this end, we augment the type system of a subset of the C
programming language with \emph{type qualifiers}: instead of having a single
pointer type, we add syntactic rules to distinguish between safe pointers that
can be dereferenced with the \texttt{*}, \texttt{->} and \texttt{[]} operators ;
and unsafe pointers that have to go through a special API provided by the
kernel, consisting most notably of the two functions \texttt{copy\_from\_user}
and \texttt{copy\_to\_user}.

\section{Systèmes de types}

La plupart des langages de programmation incorporent la notion de type, qui
permet de détecter ou d'empêcher de manipuler des données incompatibles entre
elles.

Nous avons vu dans le chapitre~\ref{cha:os} qu'au niveau du langage machine, les
seules données qu'un ordinateur manipule sont des nombres. Selon les opérations
effectuées, ils seront interprétés comme des entiers, des adresses mémoires, ou
des caractères. Pourtant il est clair que certaines opérations n'ont pas de
sens: par exemple, ajouter deux adresses, ou déréférencer le résultat d'une
division sont des comportements qu'on voudrait pouvoir empêcher.

En un mot, le but du typage est de classifier les objets et de restreindre les
opérations possibles selon la classe d'un objet : ``ne pas ajouter des pommes et
des oranges''. Le modèle qui permet cette classification est appelé
\emph{système de types} et est en général constitué d'un ensemble de
\emph{règles de typage}, comme ``un entier plus un entier égale un entier''.

Il y a deux grandes familles de systèmes de types, selon quand se fait la
vérification de types. On peut en effet l'effectuer au moment de l'exécution, ou
au contraire prévenir les erreurs à l'exécution en la faisant au moment de la
compilation (ou avant l'interprétation).

\paragraph{Typage dynamique :} dans ce cas, chaque valeur manipulée par le
programme est décorée d'une étiquette définissant comment interpréter la valeur
en question. Les règles de typage sont alors réalisées à l'exécution. Par
exemple, l'opérateur "$+$" vérifie que ces deux opérandes ont une étiquette
"entier", et construit alors une valeur obtenue en faisant l'addition des deux
valeurs, avec une étiquette "entier". Symboliquement :

\begin{verbatim}
fonction addition(x, y) {
  if ( (etiquette_type(x) != type_entier)
    || (etiquette_type(y) != type_entier)
    ) {
      erreur ("L'addition est seulement définie pour deux entiers")
  } else {
    vx = valeur(x)
    vy = valeur(y)
    v = addition_machine(vx, vy)
    return Valeur(valeur=v, type=type_entier)
  }
}
\end{verbatim}

Le langage Python~\link{python} utilise cette stratégie, qui est illustrée dans
la session interactive de la figure~\ref{fig:typage-dynamique}.

\begin{figure}
  \insertcode{typage-dynamique.pycon}
  \caption{Session Python présentant le typage dynamique}
  \label{fig:typage-dynamique}
\end{figure}

L'analyse de taintage consiste à ajouter des étiquettes au données décrivant
leur provenance, de manière à ce que chaque accès dangereux soit vérifié à
l'exécution. Cela se marie bien avec les langages typés dynamiquement, comme le
montre l'exemple célèbre du mode ``souillé'' (\emph{tainted}) de
Perl~\cite{perlCamelBook}. Il est aussi possible, mais plus délicat, d'appliquer
cette technique à des langages compilés~\cite{clause-etal-issta07,oakland10}.
% TODO [E]: à préciser

\paragraph{Typage statique :} dans ce cas on fait les vérifications à la
compilation. En quelque sorte, l'approche dynamique est pessimiste, puisqu'elle
demande de traiter très souvent le cas où les types ne sont pas corrects.
Intuitivement, dans le cas où toutes les fonctions se comportent bien, faire la
vérification est inutile. Pour vérifier ceci, on donne à chaque fonction un
contrat comme "si deux entiers sont passés, le fonction renverra un entier sans
provoquer d'erreur". Cet ensemble de contrats peut être vérifié statiquement par
le compilateur, à l'aide d'un système de types statique.

% TODO [E]: division par zéro

Ainsi la fonction "$+$" est typée \texttt{(int * int)} $→$ \texttt{int}. Les
règles permettant de vérifier le typage sont par exemple les suivantes :

\begin{itemize}
\item
  une constante entière est toujours de type \texttt{int}.
\item
  si $f$ a pour type $(t_1, …, t_n) → t$, et que chaque $e_i$ a pour type $t_i$
  ($i ∈ [1;n]$), alors $f(e_1, …, e_n)$ a pour type $t$.
\item
  si en considérant que chaque $e_i$ a pour type $t_i$ ($i ∈ [1;n]$),
  on arrive à typer le corps de $f$ et que sa valeur de retour a alors
  pour type $t$,
  alors $f$ a pour type $(t_1, …, t_n) → t$.
\end{itemize}

Cet ensemble de règles, une fois formalisé et implanté, est en général assez
efficace pour vérifier la correction d'un programme.
% TODO [E] ?

\paragraph{Typage fort ou faible :}

Si un système de types statique permet d'éliminer totalement la nécessité de
réaliser des tests de typage, on dit qu'il est \emph{fort}. Mais ce n'est que
rarement le cas. En effet, il peut y avoir des constructions au sein du langage
qui permettent de contourner le système de types, comme un opérateur de
transtypage\ref{fig:javacast}. À l'exécution, une erreur de types est levée :

\begin{figure}
  \insertcode{cast.java}
  \caption{Transtypage en Java}
  \label{fig:javacast}
\end{figure}

\begin{Verbatim}
Exception in thread "main" java.lang.ClassCastException:
    java.lang.Integer cannot be cast to java.lang.Float
        at Cast.main(Cast.java:5)
\end{Verbatim}

\paragraph{Polymorphisme :} parfois, il est trop restrictif de donner un unique
contrat à une fonction. Quel doit être le type d'une fonction ajoutant un
élément à une liste ?

En première approximation, on peut imaginer fournir une version du code par type
de données à manipuler. C'est la solution retenue par les premières versions du
langage Pascal, ce qui rendait très difficile l'écriture de
bibliothèques\cite{PascalNoEscape}.

Une amélioration peut être de générer à la compilation une fonction qui aurait
le même code mais des annotations de typage différentes : on obtient alors
plusieurs copies de la fonction qui sont chacunes spécialisées pour un type en
particulier. C'est la solution retenue par le langage C++ et son système de
\emph{templates}.

En fait, si à l'exécution toutes les données sont représentées de manière
uniforme, il n'est pas nécessaire de créer plusieurs copies de la fonction.

\begin{figure}
  \insertcode{listappend.ml}
  \caption{Fonction de concaténation de listes en OCaml.}
  \label{fig:listappend}
\end{figure}

Par exemple, la fonction de la figure~\ref{fig:listappend} n'opère que sur la
structure du type liste (en utilisant ses constructeurs \texttt{{[}{]}} et
\listcons ainsi que le filtrage) : les éléments de \texttt{lx} et \texttt{ly} ne
sont pas manipulés à part pour les transférer dans le résultat.

% TODO [E] tu parles uniquement du polymorphisme paramétrique : attention il y a
% la classe de polymorphisme et son implantation qui peuvent diverger

Elle peut donc être typée avec n'importe quelle type 
$a~\textrm{list} → a~\textrm{list} → a~\textrm{list}$. Pour introduire cette
généricité, on modifie le système de types pour transformer :

\[ ∀ a. \textrm{append} :
   a~\textrm{list} → a~\textrm{list} → a~\textrm{list}
\]

en :

\[ \textrm{append} : ∀ a.
   a~\textrm{list} → a~\textrm{list} → a~\textrm{list}
\]

% TODO [E] que veux tu dire ?

% TODO [E] où la variable de types peut prendre n'importe quelle valeur de type

Au lieu d'associer à chaque expression un type, dans certains cas on lui associe
un schéma de types, instanciable en un type concret. En quelque sorte, on
déplace la vérification du site d'appel au site de définition. Cette technique a
été décrite en premier dans~\cite{Milner78}.
Un tour d'horizon de différents types statiques, avec en particulier du
polymorphisme, on pourra se référer à~\cite{TAPL}.

\paragraph{Expressivité :} les systèmes de types les plus simples expriment des
contrats esssentiellement liés à la sûreté d'exécution, pour ne pas utiliser des
valeurs de types incompatibles entre eux. Mais il est possible d'étendre le
langage avec des annotations plus riches : par exemple en vérifiant statiquement
que des listes ne sont pas vides\cite{lightweight-static-capabilities}, ou dans
le domaine de la sécurité, d'empêcher des fuites d'information~\cite{LZ06a}.

\paragraph{Qualificateurs de types :} Dans le cas particulier des vulnérabilités
liées à une mauvaise utilisation de la mémoire, les développeurs du noyau Linux
ont ajouté un système d'annotations au code source. Un pointeur peut être décoré
d'une annotation \texttt{\_\_kernel} ou \texttt{\_\_user} selon s'il est sûr ou
pas. Celle-ci sont ignorées par le compilateur, mais un outil d'analyse statique
ad-hoc nommé Sparse~\link{sparse}~\cite{TorvaldsSparse} % TODO citer un seul
peut être utilisé pour détecter les cas les plus simples d'erreurs.
Il demande aussi au programmeur d'ajouter beaucoup d'annotations dans le
programme. % TODO quantifier

Ce système d'annotations sur les types a été formalisé sous le nom de
\emph{qualificateurs de types} : chaque type peut être décoré d'un ensemble de
qualificateurs (à la manière de \texttt{const}), et des règles de typage
permettent d'établir des propriétés sur le programme. Ces analyses ont été
implantée dans l'outil CQual~\cite{toplas-quals}. Ce système peut servir à
inférer les annotations \texttt{const}~\cite{pldi99}, à l'analyse de souillure
pour les chaîne de format~\cite{usenix01} % TODO citer Newsham sur le format ?
et des propriétés dépendantes du flot de contrôle, comme des invariants sur les
verrous~\cite{pldi02}, à rapprocher du concept de \emph{typestates}
~\cite{tse12-typestate}. Il a également été appliqué à la classe de
vulnérabilités sur les pointeurs utilisateurs dont il est ici
l'objet~\cite{cquk-usenix04}. Puisqu'elle consiste à ajouter un qualificateur à
chaque étage de type, cette approche est plus générique mais plus complexe que
la nôtre. % TODO bof

% TODO [E] préciser pourquoi

\section{Langages d'implantation}

% TODO C

Le langage C~\cite{KandR} présente des difficultés pour être analysé
automatiquement.

\begin{itemize}
\item les nombres flottants
(types \texttt{float}, \texttt{double} et \texttt{long double}) ont une
sémantique particulière, et il n'est pas correct d'approcher leur sémantique
par une sémantique dans $ℝ$. Afin d'être correct, il faut établir des domaines
spécifiques au flottant, comme~\cite{floatpoly}. Un tour d'horizon des
difficultés liées aux flottants est effectué dans~\cite{floatpitfalls}.

\item les pointeurs sur fonction
rendent floue la limite qui est habituellement présente entre instructions et
données. En leur présence il est impossible de faire une analyse de flot de
contrôle indépendante du flot de données. Pour pouvoir les traiter, il faut que
le domaine abstrait en question soit assez précis pour qu'un pointeur abstrait
se concrétise en un ensemble réduit de fonctions. Dans le cas où le domaine ne
peut pas borner l'ensemble des fonctions possibles et renvoie $\top$, l'analyse
ne peut pas continuer.

\item l'allocation dynamique
de données, présente dans le langage C par le biais des fonctions
\texttt{malloc} et \texttt{free}, modifie le modèle mémoire
nécessaire. Sans celle-ci, l'ensemble des zones mémoire possibles peut être
décrit statiquement : ce sont les noms de variable. Ce qu'introduit
\texttt{malloc} au langage, c'est une zone mémoire qui n'a pas de nom, et sur
laquelle on n'a qu'un pointeur.

% TODO, bof + expand

\item le transtypage (\emph{casts}) entre entiers et pointeurs
est particulièrement délicat à traiter. Dans les modèles abstraits, les
pointeurs sur données ou sur fonctions n'ont pas de représentation numérique,
seulement une représentation symbolique. Même dans l'exécution concrète, la
représentation numérique d'un pointeur est lié à de nombreux choix faits par
l'environnement d'exécution (comme la randomisation de l'espace d'adressage) qui
ne peuvent pas facilement être modélisés.

\end{itemize}

% TODO Ada

% TODO ASM

% TODO Rust

\section{Langages d'analyse}

Le langage C \cite{KandR,AnsiC} a été conçu pour être une sorte d'assembleur
portable, permettant décrire du code indépendamment de l'architecture sur
laquelle il sera compilé. Historiquement, c'est il a permis de créer Unix, et
ainsi de nombreux logiciels bas niveau sont écrits en C. En particulier, il
existe des compilateurs de C vers les différents langages machine pour à peu
près toutes les architectures.

\begin{figure}
  \centering

  \begin{tikzpicture}
[>=latex]
  \node (me) {Langage intermédiaire};
  \node[ above left of=me, xshift=-4cm] (fe1) {Langage source 1};
  \node[ below left of=me, xshift=-4cm] (fe2) {Langage source 2};
  \node[above right of=me, xshift=4cm] (be1) {Langage destination 1};
  \node[below right of=me, xshift=4cm] (be2) {Langage destination 2};

  \draw[->] (fe1) -- (me);
  \draw[->] (fe2) -- (me);

  \draw[->] (me) -- (be1);
  \draw[->] (me) -- (be2);

  \path (me) edge [loop below] node {Optimisations, analyses} (me);

  \coordinate (legbase) at ($ (me) + (0,-2cm) $);

  \draw[bigbrace] ($ (me.south east |- legbase) + (-2mm, 0) $)
               to node [yshift=-5mm] {Middle-end}
                  ($ (me.south west |- legbase) + (2mm, 0) $);
  \draw[bigbrace] (me.south west |- legbase)
               to node [yshift=-5mm]{Front-end}
                  (fe2.south west |- legbase);
  \draw[bigbrace] (be2.south east |- legbase)
               to node [yshift=-5mm]{Back-end}
                  (me.south east |- legbase);
\end{tikzpicture}


  \caption{Décomposition d'un compilateur : front-ends, middle-end, back-ends}
  \label{fig:middle-end}
\end{figure}

Lors de l'écriture d'un compilateur, on a besoin d'un langage intermédiaire qui
fasse l'intermédiaire entre \emph{front-end} et \emph{back-end}
(figure~\ref{fig:middle-end}). Depuis ce langage on doit pouvoir exprimer des
transformations intermédiaires sur cette représentation (analyses sémantiques,
optimisations, etc), mais aussi compiler ce langage vers un langage machine.

L'idée de prendre C comme langage intermédiaire est très séduisante, mais
malheureusement sa sémantique est trop complexe et trop peu spécifiée. Il est
donc judicieux d'utiliser un langage plus simple à cet effet. Dans de nombreux
projets, des sous-ensembles de C ont été définis pour aller dans ce sens.

Les premiers candidats sont bien entendu les représentations intermédiaires
utilisées dans les compilateurs C. Elles ont l'avantage d'accepter en plus du C
standard, les diverses extensions (GNU, Microsoft, Plan9) utilisées par la
plupart des logiciels. En particulier, le noyau Linux repose fortement sur les
extensions GNU.

\paragraph{GCC} utilise une représentation interne nommée
GIMPLE\cite{gcc-gimple}. Il s'agit d'une structure d'arbre écrite en C, reposant
sur de nombreuses macros afin de cacher les détails d'implémentation pouvant
varier entre deux versions de GCC. Cette représentation étant réputée difficile
à manipuler, le projet MELT\cite{gcc-melt} permet de générer une passe de
compilation à partir d'un dialecte de Lisp.

\paragraph{LLVM}\cite{llvm-pres} est un compilateur développé par la communauté
puis sponsorisé Apple. À la différence de GCC, sa base de code est écrite en
C++. Il utilise une représentation intermédiaire qui peut être manipulée soit
sous forme d'une structure de données C++, soit d'un fichier de code-octet
compact, soit sous forme textuelle.

\paragraph{Objective Caml}\link{ocamlManual} utilise pour sa génération de code
une représentation interne nommée Cmm, disponible dans les sources du
compilateur sous le chemin \texttt{asmcomp/cmm.mli} (il s'agit donc d'une
structure de données OCaml). Ce langage a l'avantage d'être très restreint, mais
malheureusement il n'existe pas directement de traducteur permettant de compiler
C vers Cmm.

\paragraph{C-~-}\cite{spjcmm} \link{cmm}, dont le nom est inspiré du précédent,
est un projet qui visait à unifier les langages intermédiaires utilisés par les
compilateurs. L'idée est que si un front-end peut émettre du C-~- (sous forme de
texte), il est possible d'obtenir du code machine efficace. Le compilateur
Haskell GHC utilise une représentation intermédiaire très similaire à C-~-.

Comme le problème de construire une représentation intermédiaire adaptée à une
analyse statique n'est pas nouveau, plusieurs projets ont déjà essayé d'y
apporter une solution. Puisque qu'ils sont développés en parallèle des
compilateurs, le support des extensions est en général moins important dans ces
langages.

\paragraph{CIL}\cite{NeculaCil} \link{kerneisCil} est une représentation en
OCaml d'un programme C, développée depuis 2002. Grâce à un mécanisme de
greffons, elle permet de prototyper rapidement des analyses statiques de
programmes.

\paragraph{Compcert} est un projet qui vise à produire un compilateur certifié
pour C. C'est à dire que le fait que les transformations conservent la
sémantique est prouvé. Il utilise de nombreux langages intermédiaires, dont CIL.
Pour le front-end, le langage se nomme Clight\cite{cfront}. Les passes de
middle-end, quant à elles, utilisent Cminor\cite{cminorSL}.

\paragraph{Newspeak}\cite{newspeak} est un langage intermédiaire développé par
EADS Innovation Works, et qui est spécialisé dans l'analyse de valeurs par
interprétation abstraite. Il sera décrit plus en détails dans la
section~\ref{sec:npk}.

\section{Plan de la thèse}

Cette thèse comporte trois parties.

\paragraph{La partie~\ref{part:ctx}} présente le contexte de ces travaux. Le
fonctionnement général d'un système d'exploitation y est détaillé, et les
problèmes de manipulation de pointeurs contrôlés par l'utilisateur y sont
introduits. On fait ensuite un tour d'horizon des techniques existantes
permettant de traiter ce problème.

% TODO [E] + proposition

\paragraph{La partie~\ref{part:lang}} décrit notre solution : \langname, un
langage impératif. On y décrit sa syntaxe, sa sémantique ainsi qu'un système de
types statiques. On l'étend ensuite pour capturer les problèmes d'adressage
mémoire présents dans les systèmes d'exploitation en ajoutant des pointeurs
contrôlés par l'utilisateur. Le système de types est également étendu. Pour
chacune de ces variantes, on établit la propriété de sûreté de typage reliant la
sémantique dynamique aux types statiques.

% TODO [E] dire un mot de Newspeak dès l'intro ou l'état de l'art

\paragraph{La partie~\ref{part:xp}} documente la démarche expérimentale associée
à ces travaux. L'implantation du système de types est décrite, afin que la
manière de transformer automatiquement du code C en \langname. Un cas d'étude
est déroulé, consistant d'un bug ayant touché le noyau Linux. Il est démontré
que le système de type capture précisément ce genre d'erreur de programmation.
Enfin, les possibilités d'extension tant théoriques qu'expérimentales sont
présentées.

%  - omniprésence
%  - complexité
%  - vérification
%  - taux de bugs après relecture
%  - vérification automatique obligatoire
%  - historique des techniques
%  - analyse non intrusive
%  - méthodes de développement du noyau
%  - syntaxe, sémantique, typage

% TODO [E] dire au début que c'est tendance Unix si c'est bien le cas

% TODO [E] + besoins EADS

% vim: spelllang=fr


\part{Méthodes formelles pour la sécurité}

\chapter{Systèmes d'exploitation}

\label{cha:os}
Le système d'exploitation est le programme qui permet à un système informatique
d'exécuter d'autres programmes. Son rôle est donc capital et ses responsabilités
multiples. Dans ce chapitre, nous allons voir à quoi il sert, et comment il peut
être implanté. Pour ce faire, nous présentons ici de quoi est constitué un
sustème Intel 32 bits et ce dont on se sert pour y implanter un système
d'exploitation.

\section{Architecture physique}

Un système informatique est principalement constitué d'un processeur (ou CPU
pour \emph{Central Processing Unit}), de mémoire principale (ou RAM pour
\emph{Random Access Memory}), et de divers périphériques.

Le processeur est constitué de divers registres internes qui permettent
d'encoder l'état dans lequel il se trouve : quelle est l'instruction courante
(registre \eip), quelle est la hauteur de la pile système (registre \esp), etc.
Son fonctionnement peut être vu de la manière la plus simple qui soit comme le
cycle :

\begin{itemize}
\item
  charger depuis la mémoire la prochaine instruction
\item
  \emph{(optionnel)} charger depuis la mémoire les données référencées
  par l'instruction
\item
  effectuer l'instruction
\item
  \emph{(optionnel)} stocker en la mémoire les données modifiées
\end{itemize}

Les instructions sont constituées d'un mnémonique (indicant quelle opération
faire) et d'un ensemble d'opérandes. La signification des opérandes dépend du
mnémonique, mais en général elles permettent de désigner les sources et la
destination. Les opérandes peuvent avoir plusieurs formes : une valeur
immédiate (\$4), un nom de registre (\texttt{\%eax}) ou une référence à la
mémoire (directement : \texttt{addr} ou indirectement : \texttt{(\%ecx)}
\footnote{
  Cela consiste à interpréter le contenu du regitre \ecx comme une adresse
  mémoire.
}
). On décrit les mnémoniques les plus utilisés, permettant de compiler un cœur
de langage impératif:

\begin{itemize}

\item \texttt{mov src, dst} copie le contenu de \texttt{src} dans \texttt{dst}.

\item \texttt{add src, dst} calcule la somme des contenus de \texttt{src} et
  \texttt{dst} et place ce résultat dans \texttt{dst}.

\item \texttt{push src} place \texttt{src} sur la pile, c'est-à-dire que cette
  instruction enlève au pointeur de pile \esp la taille de \texttt{src}, puis
  place \texttt{src} à l'adresse mémoire de la nouvelle valeur \esp.

\item \texttt{pop src} réalise l'opération inverse : elle charge le contenu de
  la mémoire à l'adresse \esp dans \texttt{src} puis incrémente \esp de la
  taille correspondante.

\item \texttt{jmp addr} saute à l'adresse \texttt{addr} : c'est l'équivalent de
  \texttt{mov addr, \%eip}.

\item \texttt{call addr} sert aux appels de fonction : cela revient à
  \texttt{push \%eip} puis \texttt{jmp addr}.

\item \texttt{ret} sert à revenir d'une fonction : c'est l'équivalent de
  \texttt{pop \%eip}.

\end{itemize}

Certaines de ces instructions font référence à la pile par le biais du registre
\esp. Cette zone mémoire n'est pas gérée de manière particulière. Elle permet de
gérer la pile des appels de fonction en cours grâce à la manière dont
\texttt{jmp} et \texttt{ret} fonctionnent. Elle sert aussi à stocker les
variables locales des fonctions.

À l'aide de ces quelques instructions on peut implanter des algorithmes
impératifs. Mais pour faire quelque chose de visible, comme afficher à l'écran
ou envoyer un paquet sur le réseau, cela ne suffit pas : il faut parler au reste
du matériel.

Pour ceci, il y a deux techniques principales. D'une part, certains
périphériques sont dits \emph{memory-mapped} : ils sont associés à un espace
mémoire particulier, qui ne permet pas de stocker des informations mais de lire
ou d'écrire des données dans le périphérique. Par exemple, écrire à l'adresse
$0xB8000$ permet d'écrire des caractères à l'écran. L'autre système principal
est l'utilisation des ports d'entrée/sortie. Cela correspond à des instructions
spéciales \texttt{in \%ax, port} et \texttt{out port, \%ax} où \texttt{port}
est un numéro qui correspond à un périphérique particulier. Par exemple, en
écrivant dans le port $0x60$ on peut contrôler l'état des indicateurs lumineux
du clavier PS/2.

\section{Tâches et niveaux de privilèges}
\label{sec:taches}

\subsection*{Alternance des tâches}

Sans mécanisme particulier, le processeur exécuterait uniquement une suite
d'instruction à la fois. Pour lui permettre d'exécuter plusieurs tâches, un
système de partage du temps existe.

À des intervalles de temps réguliers, le système est programmé pour recevoir une
interruption. C'est une condition exceptionnelle (au même titre qu'une division
par zéro) qui fait sauter automatiquement le processeur dans une routine de
traitement d'interruption. À cet endroit le code peut sauvegarder les registres
et restaurer un autre ensemble de registres, ce qui permet d'exécuter plusieurs
tâches de manière entrelacée. Si l'alternance est assez rapide, cela peut donner
l'illusion que les programmes s'exécutent en parallèle. Comme l'interruption
peut survenir à tout moment, on parle de multitâche préemptif.

En plus de cet ordonnancement de processus, l'architecture Intel permet
d'affecter des niveaux de privilège à ces tâches, en restreignant le type
d'instructions exécutables, ou en donnant un accès limité à la mémoire aux
tâches de niveaux moins élevés.

Il y a 4 niveaux de privilèges (nommés aussi \emph{rings}): le \ring{0} est le
plus privilégié, le \ring{3} le moins privilégié. Dans l'exemple précédent, on
pourrait isoler l'ordonnanceur de processus en le faisant s'exécuter en \ring{0}
alors que les autres tâches seraient en \ring{3}.

\subsection*{Mémoire virtuelle}

À partir du moment où plusieurs processus s'exécutent de manière concurrente, un
problème d'isolation se pose: si un processus peut lire dans la mémoire d'un
autre, des informations peuvent fuiter; et s'il peut y écrire, il peut en
détourner l'exécution.

Le mécanisme de mémoire virtuelle permet de donner à deux tâches une vue
différente de la mémoire: c'est-à-dire que vue de tâches différentes, une
adresse contiendra une valeur différente (figure~\ref{fig:memoire-virtuelle}).

\begin{figure}[h] % fig:memoire-virtuelle {{{
\centering
\begin{tikzpicture}
  [memlayout/.style={draw
                    ,shape=rectangle
                    ,minimum height=5cm
                    ,minimum width=1cm
                    }
  ,node distance=2.5cm
  ]

  \node[memlayout]            (A) {};
  \node[memlayout,right of=A] (B) {};
  \node[memlayout,right of=B] (C) {};

  \def\putlabel#1#2{
    \path (#1.north) -- ++(0,3mm) node {#2};
  }

  \putlabel{A}{Processus 1}
  \putlabel{B}{Mémoire}
  \putlabel{C}{Processus 2}

  \def\mmap#1#2#3#4#5{
    \begin{scope}
      [mzone/.style={draw
                    ,shape=rectangle
                    ,fill=black!20
                    ,minimum height=#5 mm
                    ,minimum width=1cm
                    }
      ,corr/.style={dashed}
      ]
      \node[mzone] (L) at ($ (#1.north) - (0, #2) $) {};
      \node[mzone] (R) at ($ (#3.north) - (0, #4) $) {};

      \draw[corr] (L.north east) -- (R.north west);
      \draw[corr] (L.south east) -- (R.south west);
    \end{scope}
  }

  \mmap{A}{1.5}{B}{2.5}{4}
  \mmap{A}{4.5}{B}{4}{3}
  \mmap{B}{1}{C}{1}{2}
  \mmap{B}{3.2}{C}{2.7}{4}
\end{tikzpicture}
\caption{Mécanisme de mémoire virtuelle.}
\label{fig:memoire-virtuelle}
\end{figure} % }}}

\label{page:mem-virt}
Ce mécanisme est contrôlé par valeur du registre \crtrois: les 10 premiers bits
d'une adresse virtuelle sont un index dans le répertoire de pages qui commence à
l'adresse contenue dans \crtrois. À cet index, se trouve l'adresse d'une table
de pages. Les 10 bits suivants de l'adresse sont un index dans cette page,
donnant l'adresse d'une page de 4 kio (figure~\ref{fig:pagetables}).
Plus de détails sur l'utilisation de ce mécanisme seront donnés dans la
section~\ref{sec:linux-sys}.

\begin{figure}[h]
\begin{tikzpicture}
[ bit/.style={ draw
             , shape=rectangle
             , minimum width=4mm
             , minimum height=4mm
             , node distance=4mm
             }
, pstuff/.style={ draw
                , shape=rectangle
                , node distance=1cm
                , text width=2cm
                , text centered
                }
, >=latex
]

\node [bit] (b31) {};

% generate this ?
\node [bit, right of=b31] (b30) {}; \node [bit, right of=b30] (b29) {}; \node [bit, right of=b29] (b28) {};
\node [bit, right of=b28] (b27) {}; \node [bit, right of=b27] (b26) {}; \node [bit, right of=b26] (b25) {};
\node [bit, right of=b25] (b24) {}; \node [bit, right of=b24] (b23) {}; \node [bit, right of=b23] (b22) {};
\node [bit, right of=b22] (b21) {}; \node [bit, right of=b21] (b20) {}; \node [bit, right of=b20] (b19) {};
\node [bit, right of=b19] (b18) {}; \node [bit, right of=b18] (b17) {}; \node [bit, right of=b17] (b16) {};
\node [bit, right of=b16] (b15) {}; \node [bit, right of=b15] (b14) {}; \node [bit, right of=b14] (b13) {};
\node [bit, right of=b13] (b12) {}; \node [bit, right of=b12] (b11) {}; \node [bit, right of=b11] (b10) {};
\node [bit, right of=b10] (b9) {}; \node [bit, right of=b9] (b8) {}; \node [bit, right of=b8] (b7) {};
\node [bit, right of=b7] (b6) {}; \node [bit, right of=b6] (b5) {}; \node [bit, right of=b5] (b4) {};
\node [bit, right of=b4] (b3) {}; \node [bit, right of=b3] (b2) {}; \node [bit, right of=b2] (b1) {};
\node [bit, right of=b1] (b0) {};

\foreach \b in {0,11,12,21,22,31}
    \node [above of=b\b, node distance=5mm] {\small \b};

\newcommand{\mbrace}[3]{
    \draw [bigbrace]                      ($ (b#1.south east) + (-0.1,-0.3) $)
    -- node [midway, yshift=-2mm] (#3) {} ($ (b#2.south west) +  (0.1,-0.3) $);
} % mbrace

\mbrace{0}{11}{bpp};
\mbrace{12}{21}{bpt};
\mbrace{22}{31}{bpd};

\node[pstuff, below of=bpd] (pd) {Répertoire de pages};
\node[pstuff, below of=bpt] (pt) {Table de pages};
\node[pstuff, below of=bpp] (pp) {Page de 4 kio};

\node[left of=pd, node distance=2cm] (cr) {\crtrois};

\draw [->] (bpp) -- (pp);
\draw [->] (bpt) -- (pt);
\draw [->] (bpd) -- (pd);

\draw [->] (cr) -- (pd);
\draw [->] (pd) -- (pt);
\draw [->] (pt) -- (pp);
\draw [->] (pp) -- ++ (3,0);

\end{tikzpicture}

\caption{Implantation de la mémoire virtuelle}
\label{fig:pagetables}
\end{figure}

\section{Appels système}

Avec une telle isolation, tout le code qui est exécuté en \ring{3} a une
expressivité limitée. Il ne peut pas contenir d'instructions privilégiées comme
\asminstr{in} ou \asminstr{out}, ni faire référence à des périphériques mappés
en mémoire. C'est en effet au noyau d'accéder au matériel, et pas au code
utilisateur.

Il est donc nécessaire d'appeller une routine du noyau depuis le code
utilisateur. C'est le but des appels système. Cela consiste à coupler une
fonction du \ring{3} à une fonction du \ring{0}: en appellant la fonction non
privilégiée, le flot d'exécution se retrouve dans le noyau avec les bons
privilèges.

Bien sûr, il n'est pas possible de faire directement un \asminstr{call} puisque
cela consisterait à faire un saut vers une zone plus privilégiée. Il y a
plusieurs manières d'implanter ce mécanisme. Nous décrivons ici la technique
historique à l'aide d'interruptions.

Le processeur peut répondre à des interruptions, qui sont des évenements
extérieurs. Cela permet d'écrire du code asynchrone. Par exemple, une fois qu'un
long transfert mémoire est terminé, une interruption est reçue. D'autres
interruptions dites logicielles peuvent arriver lorsqu'une erreur se produit.
Par exemple, diviser par zéro provoque l'interruption 0, et tenter d'exécuter
une instruction privilégiée provoque l'interruption 14. On peut aussi provoquer
une interruption par une instruction \asminstr{int} dédiée.

Une table globale définit, pour chaque numéro d'interruption : quelle est la
routine à appeller pour la traiter, avec quel niveau de privilège, ainsi que le
niveau de privilège requis pour pouvoir déclencher celle-ci avec l'instruction
\asminstr{int}.

Il est donc possible de créer une interruption purement logicielle (on utilise
en général le numéro 128 soit \hexa{80}), déclenchable en \ring{3} et traitée en
\ring{0}. Les registres sont préservés, donc on peut les utiliser pour passer un
numéro d'appel système (par exemple 3 pour \texttt{read()} et 5 pour
\texttt{open()}) et leurs arguments.

\section{Le \emph{Confused Deputy Problem}}
\label{sec:secu-syscalls}

On a vu que les appels systèmes permettent aux programmes utilisateur d'accéder
au services du noyau. Ils forment donc une interface particulièrement sensible aux
problèmes de sécurité.

Comme pour toutes les interfaces, on peut être plus ou moins fin. D'un
côté, une interface pas assez fine serait trop restrictive et ne permettrait pas
d'implanter tout type de logiciel. De l'autre, une interface trop laxiste
(«écrire dans tel registre matériel») empêche toute isolation. Il faut
donc trouver la bonne granularité.

Nous allons présenter ici une difficulté liée à la manipulation de mémoire au
sein de certains types d'appels système.

Il y a deux grands types d'appels systèmes: d'une part, ceux qui renvoient un
simple nombre, comme \texttt{getpid} qui renvoie le numéro du processus
appelant.

\insertcode{appel-syscall-getpid.c}

Ici, pas de difficulté particulière: la communication entre le \ring{0} et le
\ring{3} est faite uniquement à travers les registres, comme décrit dans la
section~\ref{sec:impl-syscall}.

Mais la plupart des appels systèmes communiquent de l'information de manière
indirecte, à travers un pointeur. L'appellant alloue une zone mémoire dans son
espace d'adressage et passe un pointeur à l'appel système. Ce mécanisme est
utilisé par exemple par la fonction \texttt{gettimeofday}
(figure~\ref{fig:appel-gettimeofday}).

\begin{figure}[h]
\insertcode{appel-syscall-gettimeofday.c}
\caption{Appel de gettimeofday}
\label{fig:appel-gettimeofday}
\end{figure}

Considérons une implantation naïve de cet appel système qui écrirait
directement à l'adresse pointée. Si le pointeur fourni est dans l'espace
d'adressage du processus, on est dans le cas d'utilisation normal et l'écriture
est donc possible.

Si l'utilisateur passe un pointeur dont la valeur correspond à la mémoire
réservée au noyau, que se passe-t-il? Comme le déréférencement est fait dans le
code du noyau, il est également fait en \ring{0}, et va pouvoir être réalisé
sans erreur: l'écriture se fait et potentiellement une structure importante du
noyau est écrasée.

Un utilisateur malveillant peut donc utiliser cet appel système pour écrire à
n'importe quelle adresse dans l'espace d'adressage du noyau. Ce problème vient
du fait que l'appel système utilise les privilèges du noyau au lieu de celui qui
contrôle la valeur des paramètres sensibles. Cela s'appelle le \emph{Confused
Deputy Problem}\cite{hardy88confused}.

La bonne solution est de tester dynamiquement la valeur du pointeur: s'il pointe
en espace noyau, il faut indiquer une erreur plutôt que d'écrire. Sinon, il peut
toujours y avoir une erreur, mais au moins le noyau est protégé.

\begin{figure}
\insertcode{implem-gettime.c}
\caption{Implantation de l'appel système \texttt{gettimeofday}}
\label{fig:implem-gettime}
\end{figure}

Dans le noyau, un ensemble de fonctions permet d'effectuer des copies sûres. La
fonction \texttt{access\_ok} réalise le test décrit précédemment. Les fonctions
\texttt{copy\_from\_user} et \nolinkurl{copy_to_user} réalisent une copie de la
mémoire après avoir fait ce test. Ainsi, l'implantation correcte de l'appel
système \texttt{gettimeofday} fait appel à celle-ci
(figure~\ref{fig:implem-gettime}).

Pour préserver la sécurité du noyau, il est donc nécessaire de vérifier la
valeur de tous les pointeurs dont la valeur est contrôlée par l'utilisateur.
Cette conclusion est assez contraignante, puisqu'il existe de nombreux endroits
dans le noyau où des données proviennent de l'utilisateur. Il est donc
raisonnable de vouloir vérifier automatiquement et statiquement l'absence de
tels défauts.

%\clearpage

%\section{Architecture et assembleur Intel}

%L'implantation d'un système d'exploitation est très proche du matériel sur
%lequel il s'exécute. Pour étudier une implantation en particulier, voyons ce que
%permet le matériel lui-même.

%Dans cette section nous décrivons le fonctionnement d'un processeur utilisant
%une architecture Intel 32 bits. Les exemples de code seront écrits en syntaxe
%AT\&T, celle que comprend l'assembleur GNU.\@

%La référence pour la description de l'assembleur Intel est la documentation du
%constructeur~\cite{intelsys}; une bonne explication de l'agencement dans la
%pile peut aussi être trouvée dans~\cite{SmashingTheStack}.

%Pour faire des calculs, le processeur est composé de registres, qui sont des
%petites zones de mémoire interne, et peut accéder à la mémoire principale.

%La mémoire principale contient divers types des données:

%\begin{itemize}
%\item le code des programmes à exécuter
%\item les données à disposition des programmes
%\item la pile d'appels
%\end{itemize}

%La pile d'appels est une zone de mémoire qui est notamment utilisée pour tenir
%une trace des calculs en cours. Par exemple, c'est ici que seront stockées les
%données propres à chaque fonction appelée: paramètres, adresse de retour et
%variables locales. La pile est manipulée par un pointeur de pile (\emph{stack
%pointer}), qui est l'adresse du «haut de la pile». On peut la manipuler en
%empilant des données (les placer au niveau du pointeur de pile et déplacer celui
%si) ou dépilant des données (déplacer le pointeur de pile dans l'autre sens et
%retourner la valeur présente à cet endroit).

%L'état du processeur est défini par la valeur de ses registres, qui sont des
%petites zones de mémoire interne (quelques bits chacun). Par exemple, la valeur
%du pointeur de pile est stockée dans \esp{}. Le registre \ebp{}, couplé à \esp{}
%sert à adresser les variables locales et paramètres d'une fonction, comme ce
%sera expliqué dans la section~\ref{sec:convappel}.

%L'adresse de l'instruction courante est accessible dans le registre \eip.

%En plus de ces registres spéciaux, le processeur possède de nombreux registres
%génériques, qui peuvent être utilisés pour réaliser des calculs intermédiaires.
%Ils sont nommés \eax{}, \ebx{}, \ecx{}, \edx{}, \esi{} et \edi{}. Ils peuvent
%être utilisés pour n'importe quel type d'opération, mais certains sont
%spécialisés: par exemple il est plus efficace d'utiliser \eax{} en accumulateur,
%ou \ecx{} en compteur.

%Les calculs sont décrits sous forme d'une suite d'instructions. Chaque
%instruction est composée d'un mnémonique et d'une liste d'opérandes. Les
%mnémoniques (\asminstr{mov}, \asminstr{call}, \asminstr{sub}, etc) définissent
%un type d'opération à appliquer sur les opérandes. L'instruction peut aussi être
%précédée d'une étiquette, qui correspondra à son adresse.

%\begin{center}
  %\shorthandoff{!}
  %\begin{tikzpicture}
    %[bbr/.style={bigbrace, decoration={mirror, amplitude=2mm}}
    %,asmbox/.style={minimum height=1cm}
    %,dway/.style={dashed}
    %]

    %\node[asmbox]              (lbl) {\asminstr{lbl:}};
    %\node[asmbox,right of=lbl] (mov) {\asminstr{mov}};
    %\node[asmbox,right of=mov] (op1) {\asminstr{\$4,}};
    %\node[asmbox,right of=op1] (op2) {\asminstr{\%ebx}};

    %\draw[bbr] (lbl.south west) -- (lbl.south east);
    %\draw[bbr] (mov.south west) -- (mov.south east);
    %\draw[bbr] (op1.south west) -- (op2.south east);

    %\coordinate (ops) at ($ (op1.south west)!.5!(op2.south east) $);

    %\node[below left of=lbl, xshift=-1cm, yshift=-5mm]   (tlbl) {Étiquette};
    %\node[below of=tlbl, node distance=7mm] (tmov) {Mnémonique};
    %\node[below of=tmov, node distance=7mm] (tops) {Opérandes};

    %\draw[dway] (tlbl) -| ($ (lbl.south) + (0, -0.3) $);
    %\draw[dway] (tmov) -| ($ (mov.south) + (0, -0.3) $);
    %\draw[dway] (tops) -| ($ (ops.south) + (0, -0.3) $);

  %\end{tikzpicture}
  %\shorthandon{!}
%\end{center}

%Ces opérandes peuvent être de plusieurs types:

%\begin{itemize}
%\item
%\item
  %une opérande mémoire, c'est-à-dire le contenu de la mémoire à une
  %adresse effective. Cette adresse effective peut être exprimée de
  %plusieurs manières:
  %\begin{itemize}
  %\item
    %directement: \texttt{addr}
  %\item
    %indirectement: \texttt{(\%ecx)}. L'adresse effective est le contenu
    %du registre.
  %\item
    %«base + déplacement»: \texttt{4(\%ecx)}. L'adresse effective est
    %le contenu du registre plus le déplacement (4 ici).
  %\end{itemize}
%\end{itemize}

%En pratique il y a des modes d'adressage plus complexes, et toutes les
%combinaisons ne sont pas possibles, mais ceux-ci suffiront à décrire les
%exemples suivants:

%\section{Fonctions et conventions d'appel}
%\label{sec:convappel}

%\begin{figure} % {{{ fig:stackframe
%\begin{tikzpicture}
[ stack/.style={draw,shape=rectangle,minimum height=1cm,minimum width=5cm}
, bigbrace/.style={decorate,decoration={brace,amplitude=3mm}}
, >=latex
]

  \path (0,0) node [stack] (gloc) {Locales de g}
  -- ++(0,-1) node [stack] (gebp) {}
  -- ++(0,-1) node [stack] (gret) {}
  -- ++(0,-1) node [stack] (gpar) {Arguments de g}
  -- ++(0,-1) node [stack] (floc) {Locales de f}
  -- ++(0,-1) node [stack] (febp) {}
  -- ++(0,-1) node [stack] (fret) {}
  -- ++(0,-1) node [stack] (fpar) {Arguments de f};

  \draw [->] (gebp) -- ++(4, 0) |- (febp.5);
  \draw [->] (febp.355) -- ++(1.5, 0) -- ++(0,-4);

  \draw [dashed] (fpar.south west) -- ++ (0,-2);
  \draw [dashed] (fpar.south east) -- ++ (0,-2);
  \draw [dashed] (gloc.north east) -- ++ (0,2);
  \draw [dashed] (gloc.north west) -- ++ (0,2);

  \node at (5, 1) (ox) {};

  \node [text width=3cm] at (-0.8, 2)    {Adresses hautes (sens de \asminstr{push})};
  \node [text width=3cm] at (-0.8, -8.5) {Adresses basses (sens de \asminstr{pop})};

  \node [text width=3cm] at (-6, -3) (asmtext) {\tt
    g:  ... \\
    leave \\
    ret \\
    f:  ... \\
    pushl \$2 \\
    pushl \$4 \\
    call g \\
    addl \$8, \%esp \\
    leave \\
    ret \\
    ... \\
    call f\\
    ...
    };

  \draw [-*] (gret) -- ++ (-3, 0) |- ($ (asmtext.north west) + (2, -3.5) $);
  \draw [-*] (fret) -- ++ (-3, 0) |- ($ (asmtext.north west) + (2, -6) $);

  \draw [<-] (gebp.west) -- ++ (-0.5,0) node [xshift=-4mm] {\ebp};
  \draw [<-] (gloc.west) -- ++ (-0.5,0) node [xshift=-4mm] {\esp};

  \draw [<-] ($ (asmtext.north west) + (2, 0) $) -- ++ (0.5,0) node [xshift=3mm] {\eip};

  \draw [bigbrace]
        ($ (gloc.north east) + (3, 0) $)
     -- node [auto,midway,xshift=3mm] {Cadre de g}
        ($ (gpar.south east) + (3, 0.1) $);

  \draw [bigbrace]
        ($ (floc.north east) + (3, -0.1) $)
     -- node [auto,midway,xshift=3mm] {Cadre de f}
        ($ (fpar.south east) + (3, 0) $);

\end{tikzpicture}

%\caption[Cadres de pile]{ Cadres de pile. }
%\label{fig:stackframe}
%\end{figure} % }}}

%Dans le langage d'assemblage, il n'y a pas de notion de fonction; mais
%\asminstr{call} et \asminstr{ret} permettent de sauvegarder et de restaurer une
%adresse de retour, ce qui permet de faire un saut et revenir à l'adresse
%initiale. Ce système permet déjà de créer des procédures, c'est-à-dire des
%fonctions sans arguments ni valeur de retour.

%Pour gérer ceux-ci, il faut que les deux morceaux (appelant et appelé) se
%mettent d'accord sur une convention d'appel commune. La convention utilisée sous
%GNU/Linux est appelée \emph{cdecl} et possède les caractéristiques suivantes:

%\begin{itemize}
%\item la valeur de retour d'une fonction est stockée dans \eax{}
%\item \eax{}, \ecx{} et \edx{} peuvent être écrasés sans avoir à les sauvegarder
%\item les arguments sont placés sur la pile (et enlevés) par l'appelant. Les
  %paramètres sont empilés de droite à gauche.
%\end{itemize}

%Pour accéder à ses paramètres, une fonction peut donc utiliser un adressage
%relatif à \esp. Cela peut fonctionner, mais cela complique les choses si elle
%contient aussi des variables locales. En effet, les variables locales sont
%placées sur la pile, au dessus des (c'est-à-dire, empilées après) paramètres,
%augmentant le décalage.

%Pour simplifier, la pile est organisée en cadres logiques: chaque cadre
%correspond à un niveau dans la pile d'appels de fonctions. Si $f$ appelle $g$,
%qui appelle $h$, il y aura dans l'ordre sur la pile le cadre de $f$, celui de
%$g$ puis celui de $h$.

%Ces cadres sont chainés à l'aide du registre \ebp{}: à tout moment, \ebp{}
%contient l'adresse du cadre de l'appelant.

%Prenons exemple sur la figure~\ref{fig:stackframe}: pour appeler
%\verb!g(4,2)!, \texttt{f} empile les arguments de droite à gauche.
%L'instruction \asminstr{call g} empile l'adresse de l'instruction suivante sur
%la pile puis saute au début de \texttt{g}.

%Au début de la fonction, les trois instructions permettent de sauvegarder
%l'ancienne valeur de \ebp{}, faire pointer \ebp{} à une endroit fixe dans le
%cadre de pile, puis allouer 8 octets de mémoire pour les variables locales.

%Dans le corps de la fonction \texttt{g}, on peut donc se référer aux variables
%locales par \texttt{-4(\%ebp)}, \texttt{-8(\%ebp)}, etc, et aux arguments par
  %\texttt{8(\%ebp)}, \texttt{12(\%ebp)}, etc.

%À la fin de la fonction, l'instruction \asminstr{leave} est équivalente à
%\texttt{mov \%ebp, \%esp} puis \texttt{pop \%ebp} et permet de défaire le cadre
%de pile, laissant l'adresse de retour en haut de pile. Le \asminstr{ret} final
%la dépile et y saute.


% vim: spelllang=fr


\chapter{Typage}

\input{chap/03-typage.tex}

\chapter{État de l'art}

L'analyse statique de programmes est un sujet de recherche actif depuis
l'apparition de la science informatique.

\section{Méthodes syntaxiques}

L'analyse la plus simple consiste à traiter un programme comme du texte, et à y
rechercher des motifs dangereux. Ainsi, utiliser des outils comme \texttt{grep}
permet parfois de trouver un grand nombre de vulnérabilités~\cite{SpenderGrep}.

On peut continuer cette approche en recherchant des motifs mais en étant
sensible à la syntaxe et au flot de contrôle du programme. Cette notion de
\emph{semantic grep} est présente dans l'outil Coccinelle
\cite{coccinelle09,coccinelle11} \todo{lire coccinelle09} : on peut définir des
\emph{patches sémantiques} pour détecter ou modifier des constructions
particulières.

\section{Qualificateurs de types}

Dans le cas particulier des vulnérabilités liées à une mauvaise utilisation de
la mémoire, les développeurs du noyau Linux ont ajouté un système d'annotations
au code source. Un pointeur peut être décoré d'une annotation
\texttt{\_\_kernel} ou \texttt{\_\_user} selon s'il est sûr ou pas. Celle-ci
sont ignorées par le compilateur, mais un outil d'analyse statique ad-hoc nommé
Sparse \link{sparse} peut être utilisé pour détecter les cas les plus simples
d'erreurs.

Ce système d'annotations sur les types a été formalisé sous le nom de
\emph{qualificateurs de types} : chaque type peut être décoré d'un ensemble de
qualificateurs (à la manière de \texttt{const}), et des règles de typage
permettent d'établir des propriétés sur le programme. Ces analyses ont été
implantée dans l'outil CQual
\cite{pldi99,usenix01,pldi02,cquk-usenix04,toplas-quals}. \todo{lister les
applications}.

\section{Interprétation abstraite}

L'interprétation abstraite est une technique d'analyse générique qui permet de
simuler statiquement tous les comportements d'un programme Cousot
\cite{Cousot77,Cousot92-1}. Un exemple d'application est de calculer les bornes
de variations des variables pour s'assurer qu'aucun débordement de tableau n'est
possible. Cette technique est très puissante mais possède plusieurs
inconvénients. D'une part, pour réaliser une analyse interprocédurale il faut
partir d'un point en particulier du programme (comme la fonction \texttt{main}).
Cette hypothèse n'est pas facilement satisfaite dans un noyau de système
d'exploitation, qui possède de nombreux points d'entrée. D'autre part,
il est très difficile de faire passer à l'échelle un interpréteur abstrait
\cite{AstreeScale,coverityBillion}.

\begin{center}\rule{3in}{0.4pt}\end{center}

\begin{figure}%{{{
\centering
\begin{tikzpicture}
\node at (0,0) (lbase) {};
\path (lbase) +(0,0)  node (lb) {}
              +(-1,1) node (lm) {}
              +(0,1)  node (lz) {}
              +(1,1)  node (lp) {}
              +(0,2)  node (lt) {}
;
\draw (lb) -- (lm) -- (lt) -- (lz) -- (lb) -- (lp) -- (lt);

\node at (4,-0.5) (mbase) {};
\path (mbase) +(0,0)  node (mb)  {}
              +(0,1)  node (mz)  {}
              +(-1,2) node (mle) {}
              +(1,2)  node (mge) {}
              +(0,3)  node (mt)  {}
;
\draw (mb) -- (mz) -- (mle) -- (mt) -- (mge) -- (mz);

\node at (8,-0.5) (rbase) {};
\path (rbase) +(0,0)  node (rb)  {}
              +(-1,1) node (rlt) {}
              +(0,1)  node (rz)  {}
              +(1,1)  node (rgt) {}
              +(-1,2) node (rle) {}
              +(0,2)  node (rnz) {}
              +(1,2)  node (rge) {}
              +(0,3)  node (rt)  {}
;
\draw (rb) -- (rlt) -- (rle) -- (rt) -- (rge) -- (rgt) -- (rb) -- (rz) -- (rge);
\draw (rlt) -- (rnz) -- (rgt);
\draw (rz) -- (rle);
\draw (rnz) -- (rt);

\foreach \x in {lm,lz,lp,lb,lt,mb,mle,mz,mge,mt,rb,rlt,rz,rgt,rle,rnz,rge,rt}
  {\draw (\x) circle (2pt);}



\draw[dotted] (lb) -- (mb);
\draw[dotted] (lt) -- (mt);
\draw[dotted] (lm) -- (mle);
\draw[dotted] (lp) -- (mge);
\draw[dotted] (lz) -- (mz);

\draw[dotted] (mb) -- (rb);
\draw[dotted] (mt) -- (rt);
\draw[dotted] (mle) to [bend left=30] (rle);
\draw[dotted] (mge) to [bend right=30] (rge);
\draw[dotted] (mz)  to [bend right=30] (rz);

\end{tikzpicture}

\caption{Connexions de Galois}

\end{figure}%}}}

Les domaines les plus simples ne capturent aucune relation entre variables. Ce
sont des domaines non relationnels. On peut en citer quelques uns.

\paragraph{Le domaine des signes} capture uniquement le signe des variables
(figure~\ref{fig:dom-sig}).

\begin{figure}%{{{
\centering
\begin{tikzpicture}
\node at (1, 2) (t) {$\top$};
\node at (0, 1) (p) {$+$};
\node at (2, 1) (m) {$-$};
\node at (1, 0) (b) {$\bot$};
\draw (t) -- (p) -- (b);
\draw (t) -- (m) -- (b);
\end{tikzpicture}

\begin{align*}
\gamma~(+) = & \mathbb{R}^+ \\
\gamma~(-) = & \mathbb{R}^-
\end{align*}

\caption{Domaine des signes}
\label{fig:dom-sig}
\end{figure}%}}}

\paragraph{Le domaine des intervalles} retient les bornes de variations
extremales des variables (figure~\ref{fig:dom-intvl}).

\begin{figure} % {{{
\centering
\begin{tikzpicture}
\node at (0,0)   (b)  {$\bot$};

\node at (-4,2) (cm2)   {$\{-2\}$};
\node at (-2,2) (cm1)   {$\{-1\}$};
\node at  (0,2)  (c0)   {$\{0\}$};
\node at  (2,2)  (c1)   {$\{1\}$};
\node at  (4,2)  (c2)   {$\{2\}$};

\node at (-3,3) (rm2m1) {[-2;-1]};
\node at (-1,3) (rm10)  {[-1;0]};
\node at  (1,3) (r01)   {[0;1]};
\node at  (3,3) (r12)   {[1;2]};

\node at (-2,4) (rm20)  {[-2;0]};
\node at  (0,4) (rm11)  {[-1;1]};
\node at  (2,4) (r02)   {[0;2]};

\node at (-1,5) (rm21)  {[-2;1]};
\node at  (1,5) (rm12)  {[-1;2]};

\node at  (0,6) (rm22)  {[-2;2]};

\node at  (0,9) (t)  {$\top$};

\draw (b) -- (cm2) -- (rm2m1) -- (rm20) -- (rm21) -- (rm22);
\draw (b) -- (cm1) -- (rm10) -- (rm11) -- (rm12);
\draw (b) -- (c0) -- (r01) -- (r02);
\draw (b) -- (c1) -- (r12);
\draw (b) -- (c2);

\draw (c2) -- (r12) -- (r02) -- (rm12) -- (rm22);
\draw (c1) -- (r01) -- (rm11) -- (rm21);
\draw (c0) -- (rm10) -- (rm20);
\draw (cm1) -- (rm2m1);

\draw[loosely dashed] (cm2)   to [bend left=30] (t);
\draw[loosely dashed] (rm2m1) to [bend left=30] (t);
\draw[loosely dashed] (rm20)  to [bend left=30] (t);
\draw[loosely dashed] (rm21)  to [bend left=30] (t);

\draw[loosely dashed] (rm22)  to (t);

\draw[loosely dashed] (rm12)   to [bend right=30] (t);
\draw[loosely dashed] (r02)   to [bend right=30] (t);
\draw[loosely dashed] (r12)   to [bend right=30] (t);
\draw[loosely dashed] (c2)    to [bend right=30] (t);

\draw[loosely dashed] (cm2) -- +(-1,0);
\draw[loosely dashed] (c2)  -- +(1,0);

\end{tikzpicture}

\caption{Domaine des intervalles}
\label{fig:dom-intvl}
\end{figure} % }}}

Lorsque plusieurs variables sont analysées en même temps, utiliser de tels
domaines revient à considérer un produit cartésien d'ensembles :

\begin{center}
  \begin{tikzpicture}
  \draw[->] (0,-1) -- (0,3);
  \draw[->] (-1,0) -- (4,0);

  \draw[fill=red!30] (1,1) rectangle (3,2);

  \draw[dashed] (3,1) -- (3,0);
  \draw[dashed] (1,1) -- (1,0);

  \draw[dashed] (1,2) -- (0,2);
  \draw[dashed] (1,1) -- (0,1);

  \draw[line width=3pt] (0,1) -- (0,2);
  \draw[line width=3pt] (3,0) -- (1,0);

  \end{tikzpicture}
\end{center}

Cela revient à oublier les relations entre les variables. Des domaines abstraits
plus précis permettent de retenir celles-ci. Pour ce faire, il faut modéliser
l'ensemble des valeurs des variables comme un tout.

\paragraph{Le domaines des polyèdres} est historiquement l'un des premiers
domaines relationnels. Il permet de retenir tous les invariants affines entre
fonctions (figure~\ref{fig:dom-poly}).

\begin{figure}%{{{
  \centering

  \begin{tikzpicture}[scale=0.5]
  \draw[->] (0,-1) -- (0,6);
  \draw[->] (-1,0) -- (8,0);

  \draw[fill=red!30] (1,2) -- (2,4) -- (4,5) -- (6,5) -- (7,3) -- (5,1) -- cycle;

  \end{tikzpicture}

  \[
  \vec{x} \in \mathcal{P}(A,B) \iff A \vec{x} \le B
  \]

\caption{Domaine des polyèdres}
\label{fig:dom-poly}

\end{figure}%}}}

\paragraph{Le domaine des zones} permet de représenter des relations affines de
forme $v_i - v_j \le c$ (figure~\ref{fig:dom-zones}).

\begin{figure}%{{{
  \centering
  \begin{tikzpicture}
  \draw[->] (0,-1) -- (0,5);
  \draw[->] (-1,0) -- (5,0);

  \draw[fill=red!30] (1,2) -- (3,4) -- (4,4) -- (4,3) -- (2,1) -- (1,1) -- cycle;

  \end{tikzpicture}

  \caption{Domaine des zones}
  \label{fig:dom-zones}
\end{figure}%}}}

\paragraph{Le domaine des octogones} est un compromis entre les polyèdres et les
zones. Il permet de représenter les relations $\pm v_i \pm v_j \le c$
(figure~\ref{fig:dom-octo}).

\begin{figure}%{{{
  \centering
  \begin{tikzpicture}
  \draw[->] (0,-1) -- (0,5);
  \draw[->] (-1,0) -- (5,0);

  \draw[fill=red!30] (1,1.5) -- (1,2) -- (3,4) -- (3.5,4) -- (4,3.5) -- (4,3) -- (2,1) -- (1.5,1) -- cycle;

  \end{tikzpicture}

  \caption{Domaine des octaèdres}
  \label{fig:dom-octo}
\end{figure}%}}}

\section{Typage}

L'approche par typage, plus légère, est séduisante. Pour les différents enjeux
des systèmes de types statiques, on pourra se référer à~\cite{TAPL}. Il est
possible d'encoder ce genre de propriétés dans un sytème de types, cf.
\cite{lightweight-static-capabilities} et \cite{LZ06a}.

\section{Logique de Hoare}

On peut aller plus loin que les simples types et utiliser un langage de contrats:
chaque fonction est annotée par des pré- et post-conditions sur la
mémoire\cite{cssv}\todo{Hoare}.

\section{Analyse dynamique}

Du côté de l'analyse dynamique, \cite{oakland10}\todo{et Perl ?}.

\section{Analyse de flot}

Ce que nous voulons vérifier peut être vue comme une propriété de flot. Un
\emph{survey} des problèmes et techniques existantes peut être trouvé
dans~\cite{sm-jsac03}.

\wip{}

Interprétation Abstraite :
    widening \cite{granger},
    CGS \cite{cgs},
    Astrée : presentation\cite{Astree04,Astree05},

\section{Divers}

Divers : Taint sequences \cite{mdv10},

Frama-C ?

CCurred ?


\part{Typage statique de langages impératifs}

\chapter{Un premier système de types}

Dans ce chapitre nous présentons un langage impératif permettant de modéliser
C. Nous donnons sa syntaxe, une sémantique opérationelle ainsi qu'un
système de types permettant d'obtenir plus de garanties que le système de typse
original de C tel que décrit dans \cite{AnsiC}.

Il permet le polymorphisme sur les types pointeurs, permettant par exemple de
typer :

\begin{mathpar}
⊢ \text{memcpy} : ∀ a . (a^*, a^*, \text{size\_t}) \rightarrow \tVoid
\end{mathpar}

La traduction depuis C sera explicitée dans le chapitre~\ref{cha:implem}.

\section{Syntaxe}

La grammaire suivante définit un langage impératif. On suppose qu'on peut
compiler un programme écrit en C vers ce langage.

Un programme est un triplet $P = (\vec{f}, \vec{x}, b)$ constitué d'un ensemble
de fonctions, d'un ensemble de variables et d'un bloc d'instructions. Ce bloc
sera exécuté au lancement du programme ; il peut par exemple contenir le code
d'initialisation des variables globales et l'appel à la fonction principale.

\gramlr{Programmes}{
\begin{align*}
P \gramisa & (\vec{f}, \vec{v}, b) & \textrm{Fonctions, globales, initialiseur}
\end{align*}
}

\gramlr{Expressions}{
\begin{align*}
e  \gramisa & lv     & \textrm{Left-value}
\\ \gramor  & op~e   & \textrm{Opération unaire}
\\ \gramor  & e~op~e & \textrm{Opération binaire}
\\ \gramor  & cst    & \textrm{Constante}
\\ \gramor  & \& lv  & \textrm{Pointeur sur donnée}
\\ \gramor  & \& f   & \textrm{Pointeur sur fonction}
\end{align*}
}

\gramlr{Left-values}{
\begin{align*}
lv \gramisa & var      & \textrm{Variable}
\\ \gramor  & lv.champ & \textrm{Accès à un champ}
\\ \gramor  & lv[e]    & \textrm{Accès à un tableau}
\\ \gramor  & *e       & \textrm{Déréférencement}
\end{align*}
}

\gramlr{Blocs}{
\begin{align*}
b  \gramisa & i ; b & \textrm{Séquence}
\\ \gramor  & ε     & \textrm{Bloc vide}
\end{align*}
}

\gramlr{Instructions}{
\begin{align*}
i  \gramisa & lv \leftarrow e             & \textrm{Affectation}
\\ \gramor  & lv \leftarrow funexp (args) & \textrm{Appel de fonction}
\\ \gramor  & funexp (args)               & \textrm{Appel de procédure}
\\ \gramor  & \uparrow nom \{ b \}        & \textrm{Déclaration}
\\ \gramor  & if (e) \{ b \} else \{ b \} & \textrm{Alternative}
\\ \gramor  & \{ b \} label:              & \textrm{Nommage de bloc}
\\ \gramor  & \npkGoto{label}             & \textrm{Saut en avant}
\\ \gramor  & forever \{ b\}              & \textrm{Boucle infinie}
\\ \gramor  & return e                    & \textrm{Retour de fonction}
\end{align*}
}

\todo{Fonctions}


\todo{expliquer pourquoi un while expr ne suffit pas}

La différence principale avec C est que les expressions sont sans effets de
bords, donc que leur évaluation peut être faite sans modifier l'environnement.

\section{Sémantique}

Dans cette section, nous définissons une sémantique pour ce langage impératif ;
elle pourra servir à l'implantation d'un interpréteur et à raisonner de manière
formelle sur les programme. Mathématiquement, cela consiste en la définition
d'une relation de transition $\rightarrow$ entre états de l'interpréteur.

Un état est constitué d'une part d'un point de contrôle dans le programme
(section~\ref{sec:cfg}, et d'autre part de l'état $σ$ de la mémoire
(section~\ref{sec:sigma}).

\subsection{Graphe de flot de contrôle}
\label{sec:cfg}

Dans la syntaxe ci-dessus, on peut classifier les instructions en deux familles:
celles qui définissent le flot de contrôle (\npkIf{$\cdot$}{$\cdot$}{$\cdot$},
\npkDoWith{$\cdot$}{$\cdot$}, \npkGoto{$\cdot$}, \npkWhile{$\cdot$}) et celles
qui définissent le flot de données. Une première transformation va transformer
chaque fonction en son graphe de flot de contrôle, défini comme suit :

\begin{itemize}
\item
  les nœuds sont des points de contrôle, qui représentent par exemple
  l'adresse mémoire de l'instruction qui vient d'être exécutée.
\item
  les arêtes sont soit des instructions "de données" (affectation,
  appel de fonction, déclaration), soit des conditions (ie une
  expression).
\end{itemize}

\begin{minipage}{0.5\textwidth}
\begin{Verbatim}[commandchars=\\\{\}]
\PY{k+kt}{int} \PY{n+nf}{gcd}\PY{p}{(}\PY{k+kt}{int} \PY{n}{a}\PY{p}{,} \PY{k+kt}{int} \PY{n}{b}\PY{p}{)}
\PY{p}{\PYZob{}}
    \PY{k}{if} \PY{p}{(}\PY{n}{a} \PY{o}{=}\PY{o}{=} \PY{l+m+mi}{0}\PY{p}{)} \PY{p}{\PYZob{}}
        \PY{k}{return} \PY{n}{b}\PY{p}{;}
    \PY{p}{\PYZcb{}}
    \PY{k}{while} \PY{p}{(}\PY{n}{b} \PY{o}{!}\PY{o}{=} \PY{l+m+mi}{0}\PY{p}{)} \PY{p}{\PYZob{}}
        \PY{k}{if} \PY{p}{(}\PY{n}{a} \PY{o}{\PYZgt{}} \PY{n}{b}\PY{p}{)} \PY{p}{\PYZob{}}
            \PY{n}{a} \PY{o}{=} \PY{n}{a} \PY{o}{-} \PY{n}{b}\PY{p}{;}
        \PY{p}{\PYZcb{}} \PY{k}{else} \PY{p}{\PYZob{}}
            \PY{n}{b} \PY{o}{=} \PY{n}{b} \PY{o}{-} \PY{n}{a}\PY{p}{;}
        \PY{p}{\PYZcb{}}
    \PY{p}{\PYZcb{}}
\PY{p}{\PYZcb{}}
\end{Verbatim}

\end{minipage}
\begin{minipage}{0.5\textwidth}
\begin{Verbatim}[commandchars=\\\{\}]
\PY{n}{int32} \PY{n+nf}{gcd}\PY{p}{(}\PY{n}{int32} \PY{n}{a}\PY{p}{,} \PY{n}{int32} \PY{n}{b}\PY{p}{)} \PY{p}{\PYZob{}}
    \PY{k}{if} \PY{p}{(}\PY{n}{a} \PY{o}{=}\PY{o}{=} \PY{l+m+mi}{0}\PY{p}{)} \PY{p}{\PYZob{}}
      \PY{o}{!}\PY{k}{return} \PY{o}{=} \PY{n}{b\PYZus{}int32}\PY{p}{;}
      \PY{k}{goto} \PY{n}{lbl0}\PY{p}{;}
    \PY{p}{\PYZcb{}}
    \PY{k}{while} \PY{p}{(}\PY{l+m+mi}{1}\PY{p}{)} \PY{p}{\PYZob{}}
        \PY{k}{if} \PY{p}{(}\PY{n}{b} \PY{o}{=}\PY{o}{=} \PY{l+m+mi}{0}\PY{p}{)} \PY{p}{\PYZob{}}
            \PY{k}{goto} \PY{n}{lbl1}\PY{p}{;}
        \PY{p}{\PYZcb{}}
        \PY{k}{if} \PY{p}{(}\PY{n}{a} \PY{o}{\PYZgt{}} \PY{n}{b}\PY{p}{)} \PY{p}{\PYZob{}}
             \PY{n}{a} \PY{o}{=} \PY{n}{a} \PY{o}{-} \PY{n}{b}\PY{p}{;}
        \PY{p}{\PYZcb{}} \PY{k}{else} \PY{p}{\PYZob{}}
             \PY{n}{b} \PY{o}{=} \PY{n}{b} \PY{o}{-} \PY{n}{a}\PY{p}{;}
        \PY{p}{\PYZcb{}}
    \PY{p}{\PYZcb{}}
    \PY{n+nl}{lbl1:}
    \PY{n+nl}{lbl0:}
    \PY{p}{;}
\PY{p}{\PYZcb{}}
\end{Verbatim}

\end{minipage}

\begin{tikzpicture}
\node{xxx};
\end{tikzpicture}


Intuitivement, on peut "passer" d'un état à un autre soit en passant par une
arête "condition" qui s'évalue à une valeur "vrai", soit en appliquant les
effets de bord d'une arête "instruction".

Dans la suite, on suppose qu'on a à notre disposition un ensemble de jugements :
$\langle l, instr, l' \rangle$ qui signifie qu'on peut passer du point $l$ au
point $l'$ en effectuant l'instruction $instr$.

\subsection{État mémoire}
\label{sec:sigma}

La mémoire interne de l'interpréteur est une correspondance entre l'ensemble des
adresses (infini dénombrable) et l'ensemble des valeurs. Un état mémoire $σ$ est
une fonction partielle de $Addr$ vers $Val$.

Ces valeurs peuvent être de plusieurs formes :

\gramlr{Valeurs}{
\begin{align*}
v  \gramisa  & n           & \textrm{Entier}
\\ \gramor   & f           & \textrm{Flottant}
\\ \gramor   & \cNil       & \textrm{Pointeur nul}
\\ \gramor   & \&a         & \textrm{Pointeur sur l'adresse $a$}
\\ \gramor   & \&f         & \textrm{Pointeur sur la fonction $f$}
\\ \gramor   & \top        & \textrm{Valeur non initialisée}
\end{align*}
}

\todo{Pile d'appels}

\subsection{Left values}

La mémoire est organisée en adresses, mais pourtant dans le programme cette
notion n'est pas directement visible. Les accès sont réalisés à travers des
"left values". Dans le langage C, elles correspondent aux constructions qui
peuvent se retrouver à gauche du signe "=" dans une affectation.

Une left value est donc soit une variable (locale ou globale), soit l'accès à un
indice d'une left value, soit l'accès à un champ d'une left value, soit le
déréférencement d'une left value.

\subsection{Jugements}

La sémantique concrète que nous définissons ici est constituée de jugements
logiques. Le jugement principal est une relation de transition $\rightarrow$
entre états de l'interpréteur : il sera donc noté $(l, σ) \rightarrow (l', σ')$.

Pour aider à définir ces pas, il est nécessaire d'évaluer les expressions en
fonction d'un certain état de la mémoire. On notera que dans l'état mémoire
$σ$, l'expression $e$ produit la valeur $v$ en écrivant $σ ⊢ e ⇒ v$.

Enfin, les expressions font intervenir des left-values. Puisque les états
mémoires sont définis en termes d'adresses, il faut faire la correspondance à
l'aide d'un jugement $σ ⊢ lv ⇒ a$.

\subsection{Sémantique des left-values}

\todo{Addr / Value}

Confusion : que mettre dans $σ$ : addr $\rightarrow$ value ou lv $\rightarrow$
addr ? Il faut les 2.

\begin{mathpar}
\irule{Eval-Lv-Var}{
  (v, a) ∈ σ
}{
  σ ⊢ v ⇒ a
}
\and
\irule{Eval-Lv-Deref}{
  σ ⊢ e ⇒ \&a
}{
  σ ⊢ *e ⇒ a
}
\and
\irule{Eval-Lv-Field}{
  σ ⊢ lv ⇒ a
}{
  σ ⊢ lv.f ⇒ a + f
}
\and
\irule{Eval-Lv-Array}{
  σ ⊢ lv ⇒ a \\
  σ ⊢ e ⇒ n \\
}{
  σ ⊢ lv[e] ⇒ a + n
}
\end{mathpar}

\subsection{Sémantique des expressions}

Les expressions peuvent être évaluées en valeurs. Dans le cas où l'expression
est une constante, c'est aussi le résultat.

\begin{mathpar}
\irule{Eval-Cst}{
}{
  σ ⊢ c ⇒ c
}
\end{mathpar}

Si l'expression est une left-value, on établit à quelle adresse elle correspond
et on récupère dans l'état mémoire à quelle valeur celle-ci correspond.

\begin{mathpar}
\irule{Eval-Lv}{
  σ ⊢ lv ⇒ a \\
  (a, v) ∈ σ
}{
  σ ⊢ lv ⇒ v
}
\end{mathpar}

En ce qui concerne les opérations (unaires ou binaires), on commence par évaluer
les opérandes. Le résultat est l'opération "concrète" sur les valeurs, notée
$\widehat{\textrm{op}}$. Par exemple, pour la construction syntaxique $+$, on
utilise l'addition sur les valeurs $\widehat{+}$ (c'est-à-dire l'addition
usuelle).

\begin{mathpar}
\irule{Eval-Unop}{
  σ ⊢ e ⇒ v
}{
  σ ⊢ \textrm{op}~e ⇒ \widehat{\textrm{op}}~v
}
\and
\irule{Eval-Binop}{
  σ ⊢ e_1 ⇒ v_1 \\
  σ ⊢ e_2 ⇒ v_2
}{
  σ ⊢ e_1~\textrm{op}~e_2 ⇒ v_1 \widehat{\textrm{op}}~v_2
}
\end{mathpar}

Enfin, les adresses sont aussi des valeurs. Le cas des pointeurs sur fonction
est direct puisque toutes les fonctions sont globales ; pour le cas des
pointeurs sur données on commence par déterminer l'adresse de l'objet pointé
depuis l'état mémoire.

\begin{mathpar}
\irule{Eval-AddrOfFun}{
}{
  σ ⊢ \&f ⇒ \&f
}
\and
\irule{Eval-AddrOf}{
  σ ⊢ lv ⇒ a
}{
  σ ⊢ \&lv ⇒ \&a
}
\end{mathpar}

\subsection{Sémantique des instructions}

La règle la plus simple concerne l'affectation : on peut affecter une
expressions à une left value si elles ont le même type.

\begin{mathpar}
  \irule{Instr-Assign}{
    \langle l, lv \leftarrow e, l' \rangle \\
    σ ⊢ lv ⇒ a \\
    σ ⊢ e ⇒ v
  }{
    (l, σ) \rightarrow (l', σ [ a ↦ v ])
  }
\end{mathpar}

Déclarer une variable, c'est rendre accessible dans un bloc une variable non
initialisée, qui n'est plus accessible par la suite : Si on suppose qu'on peut
traverser le bloc interne $b$ sous un $σ$ enrichi d'une nouvelle variable $x$,
on peut donc traverser l'instruction \npkDecl{x}{b}.

\begin{minipage}{0.6\textwidth}
\begin{mathpar}
  \irule{Instr-Decl}{
    \langle l, \npkDecl{x}{b}, l' \rangle \\
    \langle l_b, b, l_b' \rangle \\
    σ' = σ ⊕ \{ x \rightarrow \top \} \\
    (l_b, σ', s) \rightarrow (l_b', σ'')
  }{
    (l, σ) \rightarrow (l', σ'' \backslash x)
  }
\end{mathpar}
\end{minipage}
\begin{minipage}{0.4\textwidth}
\begin{tikzpicture}
  [cfgnode/.style={draw,shape=circle},node distance=2cm]
  \node[cfgnode] (l) {$l$};
  \node[cfgnode,right of=l, node distance=4cm] (lp) {$l'$};
  \node[cfgnode,below of=l, xshift=1cm] (lb) {$l_b$};
  \node[cfgnode,below of=lp, xshift=-1cm] (lbp) {$l_b'$};
  \draw[->] (l)  to node[auto] {$\uparrow x\{b\}$} (lp);
  \draw[->] (lb) to node[auto] {$b$} (lbp);
  \draw[->,dashed] (l) to node[auto] {$⊕ \{ x \rightarrow \top \}$}(lb);
  \draw[->,dashed,swap] (lbp) to node[auto] {$\backslash x$} (lp);
\end{tikzpicture}

\end{minipage}

\todo{fcall}

TODO pour :

\[
\irule{Instr-Fcall}{
  \langle l, lv \leftarrow fe(\vec{e}), l' \rangle \\
  σ ⊢ fe ⇒ f \\
  σ ⊢ \vec{e} ⇒ \vec{v} \\
  σ' = σ ⊕ \{args(f) = \vec{v}\} ⊕ \{ !ret \rightarrow \top \} \\
  (Entry(f), σ') \rightarrow (Exit(f), σ'') \\
  σ'' ⊢ !ret ⇒ v_{ret} \\
  σ'' ⊢ lv ⇒ a
}{
  (l, σ) \rightarrow (l', σ'' \backslash (args(f) \cup \{!ret\}) ⊕ \{ a \rightarrow v_{ret}\}, ?)
}
\]

\todo{definir $σ ⊢ fe ⇒ f$}

\subsection{Sémantique des conditions}

On utilise un encodage similaire à la déclaration. Tout d'abord, on évalue la
condition dans un contexte $σ$. Si elle s'évalue en un entier non nul, et qu'une
transition à travers le bloc $i_t$ est possible, alors on peut faire passer à
travers le "\textsc{If}".

\begin{minipage}{0.5\textwidth}
\begin{tikzpicture}
  [cfgnode/.style={draw,shape=circle}, node distance=1.5cm]
  \node[cfgnode] (l) {$l$};
  \node[cfgnode,right of=l, node distance=4.5cm] (lp) {$l'$};
  \node[cfgnode,below of=l, xshift=5mm] (lt) {$l_t$};
  \node[cfgnode,below of=lp, xshift=-5mm] (ltp) {$l_t'$};
  \node[cfgnode,above of=l, xshift=5mm] (lf) {$l_f$};
  \node[cfgnode,above of=lp, xshift=-5mm] (lfp) {$l_f'$};
  \draw[->] (l)  to node[auto] {$\iIf{e}{i_t}{i_f}$} (lp);
  \draw[->] (lt) to node[auto] {$i_t$} (ltp);
  \draw[->] (lf) to node[auto] {$i_f$} (lfp);
  \draw[->,dashed] (l) to node [auto,swap] {$≠0$} (lt);
  \draw[->,dashed] (ltp) to (lp);
  \draw[->,dashed] (l) to node [auto] {$0$} (lf);
  \draw[->,dashed] (lfp) to (lp);
\end{tikzpicture}

\end{minipage}
\begin{minipage}{0.5\textwidth}
\begin{mathpar}
\irule{If-True}{
  \langle l, \npkIf{e}{i_t}{i_f}, l' \rangle \\
  σ ⊢ e ⇒ n \\
  n ≠ 0 \\
  \langle l_i, i_t, l_i' \rangle \\
  (l_i, σ) \rightarrow (l_i', σ')
}{
  (l, σ) \rightarrow (l', σ')
}
\and
\irule{If-False}{
  \langle l, \npkIf{e}{i_t}{i_f}, l' \rangle \\
  σ ⊢ e ⇒ 0 \\
  \langle l_i, i_f, l_i' \rangle \\
  (l_i, σ) \rightarrow (l_i', σ')
}{
  (l, σ) \rightarrow (l', σ')
}
\end{mathpar}
\end{minipage}

\section{Règles de typage}

Dans cette section, on définit la notion de programme bien typé. L'analyse par
typage permet de vérifier qu'à chaque expression on peut associer un type, et ce
de manière cohérente entre plusieurs utilisations d'une variable.

Un jugement de typage est de la forme $Γ ⊢ e : τ$ et se lit "sous $Γ$, $e$ est
typable en $τ$". Un environnement de typage $Γ$ contient le contexte nécessaire
à l'analyse, c'est à dire le type des fonctions et variables du programme.

Les instructions et blocs, au contraire, n'ont pas de type. On note $Γ ⊢ i$ pour
"sous $Γ$, $i$ est bien typé", c'est à dire que ces sous expressions sont
typables et que les types sont en accord avec le flot de données (par exemple,
pour que l'instruction $lv \leftarrow e$ soit bien typée sous $Γ$, il faut que les types
de $lv$ et de $e$ puissent avoir le même type sous $Γ$).

Les types des valeurs sont :

\gramlr{Types}{
\begin{align*}
τ  \gramisa & \tInt, \tFloat, \tVoid        & \textrm{Constante}
\\ \gramor  &  a                            & \textrm{Variable}
\\ \gramor  & (τ_1, …, τ_n) \rightarrow τ_r & \textrm{Fonction}
\\ \gramor  & [τ]                           & \textrm{Tableau}
\\ \gramor  & τ*                            & \textrm{Pointeur}
\\ \gramor  & \{ f_1:τ_1
              ,    …
              , f_n:τ_n \}                  & \textrm{Structure}
\end{align*}
}

Et les environnements de typage :

\gramlr{Environnements}{
\begin{align*}
Γ  \gramisa & ε     & \textrm{Environnement vide}
\\ \gramor  & Γ,v:τ & \textrm{Ajout d'une variable}
\end{align*}
}

On va associer à chaque variable globale un type. Mais faire de même pourrait
être trop restrictif. En effet, une fonction comme memcpy peut être utilisée
pour copier des tableaux d'entiers, mais aussi de flottants. On va donc associer
un schéma de types à chaque fonction.

\gramlr{Schémas}{
\begin{align*}
σ \gramisa & ∀ \vec{a} . τ
\end{align*}
}

En associant un schéma de type $σ$ à une fonction $f$, on indique que la
fonction pourra être utilisée avec tout type $τ$ qui est une instanciation de
$σ$.

\todo{Gen, inst}

\subsection{Programme}

Au niveau global, un programme P est bien typé (noté $⊢ P$) s'il existe un
environnement $Γ = (\vec{σ}, \vec{τ})$ permettant de typer ses composantes (les fonctions, les
globales et le bloc d'initialisation).

\[
\irule{Prog}{
  (\vec{σ}, \vec{τ}) ⊢ \vec{f} : \vec{σ} \\
  (\vec{σ}, \vec{τ}) ⊢ \vec{x} : \vec{t} \\
  (\vec{σ}, \vec{τ}) ⊢ b
}{
  ⊢ (\vec{f}, \vec{x}, b)
}
\]

\subsection{Flot de contrôle}

Les règles suivantes permettent de définir les jugements $Γ ⊢ i$. De manière
générale, les instructions sont bien typées si leurs sous-instructions sont bien typées.

\begin{mathpar}
\irule{Pass}{
}{
  Γ ⊢ ε
}
\and
\irule{Seq}{
  Γ ⊢ s \\
  Γ ⊢ b
}{
  Γ ⊢ s;b
}
\and
\irule{While}{
  Γ ⊢ b
}{
  Γ ⊢ \npkWhile{b}
}
\and
\irule{Goto}{
}{
  Γ ⊢ \npkGoto{l}
}
\and
\irule{DoWith}{
  Γ ⊢ b
}{
  Γ ⊢ \npkDoWith{b}{l}
}
\end{mathpar}

Dans le cas de la conditionnelle, il est en plus nécessaire de vérifier que la
condition est un entier.

\begin{mathpar}
\irule{If}{
  Γ ⊢ e : \tInt \\
  Γ ⊢ i_t \\
  Γ ⊢ i_f
}{
  Γ ⊢ \npkIf{e}{i_t}{i_f}
}
\end{mathpar}

\subsection{Left values}

On associe à chaque left-value un type, qui est aussi le type des valeurs que
celle-ci peut contenir. Le cas des variables est direct : si un couple
(variable, type) est dans l'environnement de typage, la variable possède ce
type.

\begin{mathpar}
\irule{Lv-Var}{
  (v, τ) ∈ Γ
}{
  Γ ⊢ v : τ
}
\end{mathpar}

Si une expression a un type pointeur, en la déréférençant on obtient une valeur
du type pointé.

\begin{mathpar}
\irule{Lv-Deref}{
  Γ ⊢ e : τ*
}{
  Γ ⊢ *e : τ
}
\end{mathpar}

\todo{accès à un élément d'un type composite}

\begin{mathpar}
\irule{Lv-Field}{
  Γ ⊢ lv : τ_s \\
  (f, τ_f) ∈ τ_s
}{
  Γ ⊢ lv.f : τ_f
}
\and
\irule{Lv-Array}{
  Γ ⊢ lv : [τ] \\
  Γ ⊢ e : \tInt
}{
  Γ ⊢ lv[e] : τ
}
\end{mathpar}

\subsection{Expressions}

Les constantes ont leurs types habituels. Notons que le pointeur nul (\cNil) a
un type polymorphe.

\begin{mathpar}
\irule{Const-Int}{
}{
  Γ ⊢ n : \tInt
}
\and
\irule{Const-Float}{
}{
  Γ ⊢ f : \tFloat
}
\and
\irule{Const-Nil}{
}{
  Γ ⊢ \cNil : τ*
}
\end{mathpar}

Un certain nombre d'opérations est possible sur le type \tInt.

\begin{mathpar}
\irule{Op-Int}{
  \textrm{op} ∈ \{+,-,\times,/,\&,|,\&\&,||,⋘,⋙\} \\
  Γ ⊢ e_1 : \tInt \\
  Γ ⊢ e_2 : \tInt
}{
  Γ ⊢ e_1 \textrm{op} e_2 : \tInt
}
\end{mathpar}

De même sur \tFloat.

\begin{mathpar}
\irule{Op-Float}{
  \textrm{op} ∈ \{+.,-.,\times.,/.\} \\
  Γ ⊢ e_1 : \tFloat \\
  Γ ⊢ e_2 : \tFloat
}{
  Γ ⊢ e_1 \textrm{op}~ e_2 : \tFloat
}
\end{mathpar}

Les opérateurs de comparaison peuvent s'appliquer à deux opérandes de types
\tInt ou \tFloat et renvoient un \tInt.

\begin{mathpar}
\irule{Op-Cmp}{
  \textrm{op} ∈ \{=,≠,≤,≥,<,>\} \\
  Γ ⊢ e_1 : τ \\
  Γ ⊢ e_2 : τ \\
  τ ∈ \{ \tInt, \tFloat \}
}{
  Γ ⊢ e_1 \textrm{op} e_2 : \tInt
}
\end{mathpar}

L'opérateur unaire $-$ peut être appliqué à un \tInt ou à un \tFloat.

\begin{mathpar}
\irule{Unop-Minus}{
  τ ∈ \{ \tInt, \tFloat \} \\
  Γ ⊢ e : τ
}{
  Γ ⊢ -e : τ
}
\end{mathpar}

Les opérateurs de négation unaires, en revanche, ne s'appliquent qu'aux
entiers.

\begin{mathpar}
\irule{Unop-Not}{
  \textrm{op} ∈ \{ \textasciitilde, ! \} \\
  Γ ⊢ e : \tInt
}{
  Γ ⊢ \textrm{op}~e : \tInt
}
\end{mathpar}

On peut prendre l'adresse d'une left-value, et obtenir un pointeur vers
celle-ci. Le type du résultat est un type pointeur vers le type de base.

\begin{mathpar}
\irule{AddrOf}{
  Γ ⊢ lv : τ
}{
  Γ ⊢ \&lv : τ*
}
\end{mathpar}

La règle concernant les pointeurs sur fonction est similaire, à ceci près que
les fonctions ont des schémas de type. Un pointeur sur fonction peut donc avoir
n'importe quelle instanciation de ce schéma.

\begin{mathpar}
\irule{AddrOfFun}{
  Γ ⊢ f : σ \\
  τ ≤ σ
}{
  Γ ⊢ \&lv : τ
}
\end{mathpar}

\subsection{Fonctions}

Pour typer une fonction, on commence par ajouter ses paramètres dans
l'environnement de typage, et on type la définition de la fonction. Le type
résultant est généralisé.

\[
\irule{Fun}{
  Γ' = Γ ⊕ \{ args(f) = \vec{τ} \} \\
  Γ' ⊢ body(f) \\
  Γ' ⊢ !ret : τ_r
}{
  Γ ⊢ f : Gen(\vec{τ} \rightarrow τ_r, Γ)
}
\]

\subsection{Instructions}

\begin{mathpar}
\irule{Decl}{
  Γ ⊕ \{x : τ\} ⊢ b
}{
  Γ ⊢ \npkDecl{x}{b}
}
\and
\irule{Assign}{
  Γ ⊢ lv : τ \\
  Γ ⊢ e  : τ
}{
  Γ ⊢ lv \leftarrow e
}
\and
\irule{Fcall}{
  Γ ⊢ lv : τ_{ret} \\
  Γ ⊢ fe : σ \\
  Γ ⊢ \vec{e} : \vec{τ} \\
  (\vec{τ} \rightarrow τ_r) ≤ σ
}{
  Γ ⊢ lv \leftarrow fe(\vec{e})
}
\end{mathpar}

\section{Limitations}
\subsection{Programmes non typables}
\subsection{Incohérences}


\chapter{Système avec le qualificateur "user"}

\section{Éditions et ajouts}
\section{Passage sur le cas d'étude}

\section{Propriété d'isolation mémoire}

Le déréférencement d'un pointeur dont la valeur est contrôlée par l'utilisateur
ne peut se faire qu'à travers une fonction qui vérifie la sûreté de celui-ci.


\chapter{Système avec le qualificateur "sz"}

Dans ce chapitre, nous présentons une autre extension au système de types du
chapitre~\ref{cha:typbase}, similaire à celle du chapitre~\ref{cha:qualuser}.
Il s'agit cette fois-ci de détecter les pointeurs sur caractères (\texttt{char
*}) qui sont terminés par un caractère NUL et donc une chaîne C correcte. La
bibliothèque C propose quantité de fonctions manipulant ces chaînes et appeler
une fonction comme \texttt{strcpy} sur un pointeur quelconque est un problème de
sécurité que nous cherchons à détecter.

\section{But}
\section{Annotation de \texttt{string.h}}


\part{Expérimentation}

\chapter{Implantation}

\label{cha:implem}
\section{Langages intermédiaires}

Le langage C \cite{KandR,AnsiC} a été conçu pour être une sorte d'assembleur
portable, permettant décrire du code indépendamment de l'architecture sur
laquelle il sera compilé. Historiquement, c'est il a permis de créer Unix, et
ainsi de nombreux logiciels bas niveau sont écrits en C. En particulier, il
existe des compilateurs de C vers les différents langages machine pour à peu
près toutes les architectures.

Lors de l'écriture d'un compilateur, on a besoin d'un langage intermédiaire qui
fasse l'intermédiaire entre \emph{frontend} et \emph{backend}. Avec ce langage,
on doit pouvoir (facilement) réaliser les types d'opérations suivantes :

\begin{itemize}
\item
  compiler ce langage vers un langage machine (interface avec le
  \emph{backend}).
\item
  exprimer des transformations intermédiaires sur cette représentation
  (analyses sémantiques, optimisations, etc).
\end{itemize}

À cause du premier point, un langage comme C est très séduisant, mais
malheureusement il a de nombreux défauts qui font qu'il ne remplit pas le
deuxième. Par exemple, certaines instructions sont ambiguës, et l'ensemble des
opération n'est pas orthogonal (on peut faire la même chose de plusieurs
manières).

Pour pallier ce problème, une première solution peut être d'ajouter des
contraintes au C intermédiaire manipulé pour en obtenir un sous-langage. De
nombreux sous-ensembles ont été définis pour aller dans ce sens.

\subsubsection{Critères}

Les critères à évaluer sont les suivants :

\paragraph{Forme (textuel / langage d'implem)}

Quelle est la représentation du langage ? Est-ce une bibliothèque (si oui, dans
quel langage) a-t'il une forme textuelle (si oui, pour quels langages des
\emph{parsers} existent ils).

\paragraph{Maturité}

\paragraph{Tools}

\paragraph{Support}

Le langage a-t'il ``fait ses preuves'' ? Est-il susceptible de changer ?

\paragraph{Scope (analyse / compil)}

Pour quel type d'utilisation a-t'il été conçu ?

\paragraph{Expressivité}

Peut-on compiler ``tout C'' ?

\paragraph{Orthogonalité}

\paragraph{Typage}

\subsubsection{Langages}

\paragraph{LLVM}

LLVM\cite{llvm-pres} est un backend de compilateur développé par Apple.

\paragraph{Cmm (Haskell)}

\paragraph{Cmm (OCaml)}

\paragraph{CIL}

\paragraph{Clight}

\paragraph{Cminor}

\paragraph{Cminusminus}

\cite{spjcmm}

http://www.cminusminus.org/

\paragraph{Newspeak}

\section{Newspeak}

\section{Chaîne de compilation}

\begin{figure}
  \begin{tikzpicture}\shorthandoff{!}
\tikzstyle{file}=[draw, shape=rectangle, node distance=2.2cm, minimum
height=1cm, shade, top color=white,
    bottom color=blue!50!black!20, draw=blue!40!black!60, very thick];

\node [file] (c1) {\textcolor{black}{.c}};
\node [file, below of=c1] (c2) {\textcolor{black}{.c}};
\node [node distance=2.2cm, below of=c2] (c3) {};

\node [file, minimum height=0, node distance=5mm, above of=c3,draw] (c3b) {.adb};
\node [file, minimum height=0, node distance=5mm, below of=c3,draw] (c3s) {.ads};

\path (c3b.north west) ++(-3mm,3mm) [draw,dotted] rectangle ($(c3s.south east)+(3mm,-3mm)$);

\node [below of=c3, node distance=2.2cm] (c4){};

\node [file, right of=c1] (cc1) {\textcolor{black}{.c}};
\node [file, right of=c2] (cc2) {\textcolor{black}{.c}};
\node [node distance=2.2cm, right of=c3] (cc3) {};

\node [below of=cc3, node distance=2.2cm](cc4){};

\draw[->] (c1) -- node[above] {{\tiny \ttfamily gcc -E}} (cc1);
\draw[->] (c2) -- (cc2);

\node [file, right of=cc1] (no1) {\textcolor{black}{.no}};
\node [file, right of=cc2] (no2) {\textcolor{black}{.no}};
\node [file, right of=cc3] (no3) {\textcolor{black}{.no}};

\node [below of=no3, node distance=2.2cm](no4){};

\draw[->] (cc1) -- node[above] {{\tiny \ttfamily c2newspeak -c}} (no1);
\draw[->] (cc2) --  (no2);

\draw[->] ($ (c3b.north east)!0.5!(c3s.south east) + (3mm,0) $) -- node[above] {\tiny \ttfamily ada2newspeak -c} (no3);


\node [file, right of=no2] (npk) {\textcolor{black}{.npk}};

\node [right of=no3, node distance=2.2cm](npk2){};
\node [right of=no4, node distance=2.2cm](npk3){};

\node[draw, ellipse, right of=no2, minimum height=3cm]{};

\draw[->] (no1) -- node {} (npk);
\draw[->] (no2) -- node[above] {{\tiny \ttfamily c2newspeak}} (npk);
\draw[->] (no3) -- node {} (npk);


\node [file, right of=npk](warn) {\textcolor{black}{.warn}};

\node [right of=npk3, node distance=3cm](analyser2){};

\draw[->] (npk) -- node[above] {{\tiny \ttfamily penjili}} (warn);

\draw[->] (c4)   -- node[above, text depth=3pt] {\footnotesize Pr\'etraitement} (cc4);
\draw[->] (cc4)  -- node[above, text depth=3pt] {\footnotesize Compilation} (no4);
\draw[->] (no4)  -- node[above, text depth=3pt] {\footnotesize \'Edition de liens} (npk3);
\draw[->] (npk3) -- node[above, text depth=3pt] {\footnotesize Analyse} (analyser2);
\end{tikzpicture}

  \caption{Compilation depuis Newspeak}
  \label{fig:compil-npk}
\end{figure}

La compilation vers C est faite en trois étapes (figure~\ref{fig:compil-npk}) :
prétraitement du code source, compilation de C prétraité vers \newspeak{}, puis
compilation de \newspeak{} vers ce langage.

\subsection{Prétraitement}

\ctonewspeak{} travaillant uniquement sur du code prétraité (dans directives de
préprocesseur), la première étape consiste donc à faire passer le code par \cpp.

\subsection{Compilation (levée des ambigüités)}

Cette passe est réalisée par l'utilitaire \ctonewspeak{}. L'essentiel de la
compilation consiste à mettre à plat les définition de types, et à simplifier le
flôt de contrôle. C en effet propose de nombreuses constructions ambigües ou
redondantes.

Au contraire, \newspeak{} propose un nombre réduit de constructions. Rappelons
que le but de ce langage est de faciliter l'analyse statique : des constructions
orthogonales permettent donc d'éviter la duplication de règles sémantique, ou de
code lors de l'implémentation d'un analyseur.

Par exemple, plutôt que de fournir une boucle while, une boucle do/while et une
boucle for, \newspeak{} fournit une unique boucle \npkWhile{}. La sortie de
boucle est compilée vers un \npkGoto{}, qui est toujours un saut vers l'avant
(similaire à un "break" généralisé).

La sémantique de \newspeak{} et la traduction de C vers \newspeak{} sont
décrites dans \cite{newspeak}. En ce qui concerne l'élimination des sauts vers
l'arrière, on peut se référer à \cite{goto}.

\subsection{Annotations}

\newspeak{} a de nombreux avantages, mais pour une analyse par typage il est
trop bas niveau. Par exemple, dans le code suivant

\begin{Verbatim}[commandchars=\\\{\}]
\PY{k}{struct} \PY{n}{s} \PY{p}{\PYZob{}}
    \PY{k+kt}{int} \PY{n}{a}\PY{p}{;}
    \PY{k+kt}{int} \PY{n}{b}\PY{p}{;}
\PY{p}{\PYZcb{}}\PY{p}{;}

\PY{k+kt}{int} \PY{n+nf}{main}\PY{p}{(}\PY{k+kt}{void}\PY{p}{)}
\PY{p}{\PYZob{}}
    \PY{k}{struct} \PY{n}{s} \PY{n}{x}\PY{p}{;}
    \PY{k+kt}{int} \PY{n}{y}\PY{p}{[}\PY{l+m+mi}{10}\PY{p}{]}\PY{p}{;}
    \PY{n}{x}\PY{p}{.}\PY{n}{b} \PY{o}{=} \PY{l+m+mi}{1}\PY{p}{;}
    \PY{n}{y}\PY{p}{[}\PY{l+m+mi}{1}\PY{p}{]} \PY{o}{=} \PY{l+m+mi}{1}\PY{p}{;}
    \PY{k}{return} \PY{l+m+mi}{0}\PY{p}{;}
\PY{p}{\PYZcb{}}
\end{Verbatim}


\subsection{Implantation de l'algorithme de typage}

L'inférence de types se fait par unification, en suivant les règles de typage.

\insertcode{ex-unif-c.c}

\insertcode{ex-unif-tpk.ml}


\chapter{Étude de cas : un driver graphique}
\input{chap/09-etudedecas.tex}

\chapter{Conclusion}

\input{chap/10-conclusion.tex}

\backmatter

%\chapterstar{Liste des figures}

\listoffigures

%\chapterstar{Bibliographie}

\bibliographystyle{alpha}
\bibliography{typ}
\end{document}
